% Source: Hatcher, "Algebraic Topology", Chapter 1 "The Fundamental Group"
% (2002 Cambridge University Press edition), PDF pages 30-105.
% Transcribed page-by-page by vision subagents, then merged and restructured into
% leanblueprint conventions (semantic \label, bracketed titles, \uses dependency edges) below.
% This file is built up incrementally in chunks following the book's own subsection
% boundaries; see label_map.json for the running Hatcher-numbering <-> label map.

Algebraic topology can be roughly defined as the study of techniques for forming
algebraic images of topological spaces. Most often these algebraic images are groups,
but more elaborate structures such as rings, modules, and algebras also arise. The
mechanisms that create these images -- the `lanterns' of algebraic topology, one might
say -- are known formally as \textit{functors} and have the characteristic feature that they
form images not only of spaces but also of maps. Thus, continuous maps between
spaces are projected onto homomorphisms between their algebraic images, so
topologically related spaces have algebraically related images.

With suitably constructed lanterns one might hope to be able to form images with
enough detail to reconstruct accurately the shapes of all spaces, or at least of large
and interesting classes of spaces. This is one of the main goals of algebraic topology,
and to a surprising extent this goal is achieved.

This first chapter introduces one of the simplest and most important functors
of algebraic topology, the fundamental group, which creates an algebraic image of a
space from the loops in the space, the paths in the space starting and ending at the
same point.

\section*{The Idea of the Fundamental Group}

To get a feeling for what the fundamental group is about, let us look at a few
preliminary examples before giving the formal definitions.

Consider two linked circles $A$ and $B$ in $\R^3$, as shown in the figure. Our experience
with actual links and chains tells us that since the two circles are linked, it is impossible
to separate $B$ from $A$ by any continuous motion of $B$, such as pushing, pulling, or
twisting. We could even take $B$ to be made of rubber or stretchable string and allow
completely general continuous deformations of $B$, staying in the complement of $A$
at all times, and it would still be impossible to pull $B$ off $A$. At least that is what
intuition suggests, and the fundamental group will give a way of making this intuition
mathematically rigorous.



Instead of having $B$ link with $A$ just once, we could make it link with $A$ two or
more times. As a further variation, by assigning an orientation to $B$ we can speak of
$B$ linking $A$ a positive or a negative number of times, say positive when $B$ comes
forward through $A$ and negative for the reverse direction. Thus for each nonzero
integer $n$ we have an oriented circle $B_n$ linking $A$ $n$ times, where by `circle' we mean
a curve homeomorphic to a circle. To complete the scheme, we could let $B_0$ be a circle
not linked to $A$ at all.



Now, integers not only measure quantity, but they form a group under addition.
Can the group operation be mimicked geometrically with some sort of addition
operation on the oriented circles $B$ linking $A$? An oriented circle $B$ can be thought
of as a path traversed in time, starting and ending at the same point $x_0$, which we
can choose to be any point on the circle. Such a path starting and ending at the same
point is called a \textit{loop}. Two different loops $B$ and $B'$ both starting and ending at
the same point $x_0$ can be `added' to form a new loop $B + B'$ that travels first around
$B$, then around $B'$. For example, if $B_1$ and $B_1'$ are loops each linking $A$ once in
the positive direction, then their sum $B_1 + B_1'$ is deformable to $B_2$, linking $A$ twice.
Similarly, $B_1 + B_{-1}$ can be deformed to the loop $B_0$, unlinked from $A$. More
generally, we see that $B_m + B_n$ can be deformed to $B_{m+n}$ for arbitrary integers $m$
and $n$.



Note that in forming sums of loops we produce loops that pass through the
basepoint more than once. This is one reason why loops are defined merely as
continuous paths, which are allowed to pass through the same point many times. So
if one is thinking of a loop as something made of stretchable string, one has to give
the string the magical power of being able to pass through itself unharmed. However,
we must be sure not to allow our loops to intersect the fixed circle $A$ at any time,
otherwise we could always unlink them from $A$.

\begin{example}[The Borromean rings and the commutator $aba^{-1}b^{-1}$]
\label{ex:borromean-rings-commutator}
\group{Fundamental Group}
Next we consider a slightly more complicated sort of linking, involving three
circles forming a configuration known as the Borromean rings. The interesting
feature here is that if any one of the three circles is removed, the other two are not
linked. Let us regard one of the circles, say $C$, as a loop in the complement of the
other two, $A$ and $B$, and ask whether $C$ can be continuously deformed to unlink it
completely from $A$ and $B$, always staying in the complement of $A$ and $B$ during the
deformation.



We can redraw the picture by pulling $A$ and $B$ apart, dragging $C$ along, and then we
see $C$ winding back and forth between $A$ and $B$. Starting at the marked point of $C$
and proceeding in the direction of the arrow, we pass in sequence: (1) forward
through $A$, (2) forward through $B$, (3) backward through $A$, and (4) backward
through $B$. If we measure the linking of $C$ with $A$ and $B$ by two integers, then the
`forwards' and `backwards' cancel and both integers are zero. This reflects the fact
that $C$ is not linked with $A$ or $B$ individually.

To get a more accurate measure of how $C$ links with $A$ and $B$ together, we regard
the four parts (1)--(4) of $C$ as an ordered sequence. Taking into account the
directions in which these segments of $C$ pass through $A$ and $B$, we may deform $C$
to the sum $a+b-a-b$ of four loops. We write the third and fourth loops as the
negatives of the first two since they can be deformed to the first two but with the
opposite orientations, and the sum of two oppositely oriented loops is deformable to
a trivial loop, not linked with anything.



We would like to view the expression $a+b-a-b$ as lying in a nonabelian group, so
that it is not automatically zero. Changing to the more usual multiplicative notation
for nonabelian groups, it would be written $aba^{-1}b^{-1}$, the \textbf{commutator} of $a$ and $b$.

To shed further light on this example, suppose we modify it slightly so that the
circles $A$ and $B$ are now themselves linked. The circle $C$ can then be deformed into a
position where it again represents the composite loop $aba^{-1}b^{-1}$, where $a$ and $b$ are
loops linking $A$ and $B$. But it is apparent that $C$ can actually be unlinked completely
from $A$ and $B$ in this case. So here the product $aba^{-1}b^{-1}$ should be trivial.


\end{example}

The fundamental group of a space $X$ will be defined so that its elements are loops
in $X$ starting and ending at a fixed basepoint $x_0 \in X$, but two such loops are
regarded as determining the same element of the fundamental group if one loop can
be continuously deformed to the other within the space $X$. (All loops that occur
during deformations must also start and end at $x_0$.) In \cref{ex:borromean-rings-commutator}
above, $X$ is the complement of the circle $A$ (or of the two circles $A$ and $B$).
In \cref{sec:van-kampen-theorem} we will show:
\begin{itemize}
\item The fundamental group of the complement of a single circle $A$ is infinite cyclic
with the loop $B$ as a generator.
\item The fundamental group of the complement of two unlinked circles $A$ and $B$ is
the nonabelian free group on two generators, represented by the loops $a$ and $b$
linking $A$ and $B$. In particular, the commutator $aba^{-1}b^{-1}$ is a nontrivial element
of this group.
\item The fundamental group of the complement of two linked circles $A$ and $B$ is the
free abelian group on two generators, represented by $a$ and $b$.
\end{itemize}

As a result of these calculations, we have two ways to tell when a pair of circles $A$
and $B$ is linked: directly, by regarding one circle as an element of the fundamental
group of the complement of the other; or by distinguishing linked from unlinked
pairs via whether the fundamental group of their complement is abelian or
nonabelian. This second method is much more general: one can often show that two
spaces are not homeomorphic by showing that their fundamental groups are not
isomorphic, since it will be an easy consequence of the definition of the fundamental
group that homeomorphic spaces have isomorphic fundamental groups.

\section{1.1 Basic Constructions}
\label{sec:basic-constructions}

This first section begins with the basic definitions and constructions, and then
proceeds quickly to an important calculation, the fundamental group of the circle.

\subsection*{Paths and Homotopy}

The fundamental group will be defined in terms of loops and deformations of
loops. Sometimes it will be useful to consider more generally paths and their
deformations, so we begin with this slight extra generality.

\begin{definition}[Path]
\label{def:path}
\group{Fundamental Group}
By a \textbf{path} in a space $X$ we mean a continuous map $f : I \to X$ where $I$ is the
unit interval $[0,1]$.
\end{definition}

\begin{definition}[Homotopy of paths and homotopic paths]
\label{def:homotopy-of-paths}
\group{Fundamental Group}
\uses{def:homotopy,def:path}
A \textbf{homotopy of paths in $X$} is a family $f_t : I \to X$, $0 \le t \le 1$, such that
\begin{enumerate}
\item[(1)] the endpoints $f_t(0) = x_0$ and $f_t(1) = x_1$ are independent of $t$, and
\item[(2)] the associated map $F : I \times I \to X$ defined by $F(s,t) = f_t(s)$ is continuous.
\end{enumerate}
When two paths $f_0$ and $f_1$ are connected in this way by a homotopy $f_t$, they are
said to be \textbf{homotopic}, written $f_0 \simeq f_1$.
\end{definition}



\begin{example}[Linear homotopies in $\R^n$]
\label{ex:linear-homotopies-convex}
\dcref{ch1:1.1}
\group{Fundamental Group}
\uses{def:homotopy-of-paths}
Any two paths $f_0$ and $f_1$ in $\R^n$ having the same endpoints $x_0$ and $x_1$ are
homotopic via the homotopy $f_t(s) = (1-t)f_0(s) + t f_1(s)$. During this homotopy each
point $f_0(s)$ travels along the line segment to $f_1(s)$ at constant speed. If $f_1(s)$
happens to equal $f_0(s)$ then this segment degenerates to a point and $f_t(s) = f_0(s)$
for all $t$; this occurs in particular for $s=0$ and $s=1$, so each $f_t$ is a path from $x_0$
to $x_1$. Continuity of the homotopy $f_t$ as a map $I \times I \to \R^n$ follows from
continuity of $f_0$ and $f_1$ since vector addition and scalar multiplication are
continuous.

This construction shows more generally that for a convex subspace $X \subset \R^n$, all
paths in $X$ with given endpoints $x_0$ and $x_1$ are homotopic, since if $f_0$ and $f_1$
lie in $X$ then so does the homotopy $f_t$.
\end{example}

Before proceeding further we need to verify a technical property.

\begin{proposition}[Homotopy of paths is an equivalence relation]
\label{prop:homotopy-paths-equivalence-relation}
\dcref{ch1:1.2}
\group{Fundamental Group}
\uses{def:homotopy-of-paths}
The relation of homotopy on paths with fixed endpoints in any space is an
equivalence relation.
\end{proposition}

The equivalence class of a path $f$ under this relation is denoted $[f]$ and called the
\textbf{homotopy class} of $f$.

\begin{proof}
Reflexivity is evident since $f \simeq f$ by the constant homotopy $f_t = f$. Symmetry
is also easy since if $f_0 \simeq f_1$ via $f_t$, then $f_1 \simeq f_0$ via the inverse homotopy
$f_{1-t}$. For transitivity, if $f_0 \simeq f_1$ via $f_t$ and $f_1 \simeq g_1$ via $g_t$ with $f_1 = g_0$,
then $f_0 \simeq g_1$ via the homotopy $h_t$ that equals $f_{2t}$ for $0 \le t \le 1/2$ and $g_{2t-1}$
for $1/2 \le t \le 1$. These two definitions agree for $t=1/2$ since $f_1 = g_0$. Continuity of
the associated map $H(s,t) = h_t(s)$ comes from the fact that a function defined on
the union of two closed sets is continuous if it is continuous when restricted to each
of the closed sets separately: here $H(s,t) = F(s,2t)$ for $0 \le t \le 1/2$ and
$H(s,t) = G(s,2t-1)$ for $1/2 \le t \le 1$, where $F$ and $G$ are the maps associated to
$f_t$ and $g_t$. Since $H$ is continuous on $I \times [0,1/2]$ and on $I \times [1/2,1]$, it is
continuous on $I \times I$.
\end{proof}



\begin{definition}[Composition (product) of paths]
\label{def:path-composition}
\group{Fundamental Group}
\uses{def:path}
Given two paths $f,g : I \to X$ such that $f(1) = g(0)$, there is a \textbf{composition} or
\textbf{product path} $f \cdot g$ that traverses first $f$ and then $g$, defined by
\[
f \cdot g(s) = \begin{cases} f(2s), & 0 \le s \le 1/2 \\ g(2s-1), & 1/2 \le s \le 1. \end{cases}
\]
\end{definition}



This product operation respects homotopy classes since if $f_0 \simeq f_1$ and
$g_0 \simeq g_1$ via homotopies $f_t$ and $g_t$, and if $f_0(1) = g_0(0)$ so that $f_0 \cdot g_0$ is
defined, then $f_t \cdot g_t$ is defined and provides a homotopy $f_0 \cdot g_0 \simeq f_1 \cdot g_1$.

\begin{definition}[Loop, basepoint, and the set $\pi_1(X,x_0)$]
\label{def:loop-and-pi1-set}
\group{Fundamental Group}
\uses{def:path,def:homotopy-of-paths}
A path $f : I \to X$ with $f(0) = f(1) = x_0 \in X$ is called a \textbf{loop}, and the
common starting and ending point $x_0$ is the \textbf{basepoint}. The set of all homotopy
classes $[f]$ of loops $f : I \to X$ at the basepoint $x_0$ is denoted $\pi_1(X,x_0)$.
\end{definition}

\begin{proposition}[$\pi_1(X,x_0)$ is a group]
\label{prop:pi1-is-a-group}
\dcref{ch1:1.3}
\group{Fundamental Group}
\uses{def:loop-and-pi1-set,def:path-composition}
$\pi_1(X,x_0)$ is a group with respect to the product $[f][g] = [f \cdot g]$.
\end{proposition}

This group is called the \textbf{fundamental group} of $X$ at the basepoint $x_0$. We will
see in Chapter 4 that $\pi_1(X,x_0)$ is the first in a sequence of groups $\pi_n(X,x_0)$,
called homotopy groups.

\begin{proof}
By restricting attention to loops with a fixed basepoint $x_0$ we guarantee that the
product $f \cdot g$ of any two such loops is defined, and we have already observed that
$[f \cdot g]$ depends only on $[f]$ and $[g]$, so the product $[f][g] = [f \cdot g]$ is
well-defined. It remains to verify the three group axioms.

As a preliminary step, define a \textbf{reparametrization} of a path $f$ to be a
composition $f\varphi$ where $\varphi : I \to I$ is continuous with $\varphi(0)=0$, $\varphi(1)=1$.
Reparametrizing a path preserves its homotopy class since $f\varphi \simeq f$ via the
homotopy $f\varphi_t$ where $\varphi_t(s) = (1-t)\varphi(s) + ts$, so $\varphi_0 = \varphi$ and $\varphi_1(s) = s$; note
$(1-t)\varphi(s)+ts$ lies between $\varphi(s)$ and $s$, hence in $I$, so $f\varphi_t$ is defined.

If $f,g,h$ are paths with $f(1)=g(0)$ and $g(1)=h(0)$, then $f\cdot(g\cdot h)$ is a
reparametrization of $(f\cdot g)\cdot h$ by a piecewise linear function, hence
$(f\cdot g)\cdot h \simeq f\cdot(g\cdot h)$. Restricting to loops at $x_0$, this says the product in
$\pi_1(X,x_0)$ is associative.

Given a path $f : I \to X$, let $c$ be the constant path at $f(1)$. Then $f \cdot c$ is a
reparametrization of $f$, so $f \cdot c \simeq f$; similarly $c \cdot f \simeq f$ for $c$ the constant path
at $f(0)$. Taking $f$ to be a loop, the homotopy class of the constant path at $x_0$ is a
two-sided identity in $\pi_1(X,x_0)$.

For a path $f$ from $x_0$ to $x_1$, the \textbf{inverse path} $\bar f$ from $x_1$ back to $x_0$ is
$\bar f(s) = f(1-s)$. To see $f \cdot \bar f \simeq$ constant, use the homotopy $h_t = f_t \cdot g_t$ where
$f_t$ equals $f$ on $[0,1-t]$ and is stationary at $f(1-t)$ on $[1-t,1]$, and $g_t$ is the
inverse path of $f_t$; equivalently, describe $H(s,t)$ on the decomposition of $I\times I$
by a `V'-shaped curve, independent of $t$ below the V and independent of $s$ above
it. Since $f_0 = f$ and $f_1$ is the constant path $c$ at $x_0$, $h_t$ is a homotopy from
$f \cdot \bar f$ to $c \cdot \bar c = c$. Replacing $f$ by $\bar f$ gives $\bar f \cdot f \simeq c'$ for $c'$ the constant path
at $x_1$. Taking $f$ to be a loop at $x_0$, $[\bar f]$ is a two-sided inverse for $[f]$ in
$\pi_1(X,x_0)$.
\end{proof}

\begin{example}[$\pi_1$ of a convex set is trivial]
\label{ex:convex-set-trivial-pi1}
\dcref{ch1:1.4}
\group{Fundamental Group}
\uses{ex:linear-homotopies-convex,def:loop-and-pi1-set}
For a convex set $X$ in $\R^n$ with basepoint $x_0 \in X$ we have $\pi_1(X,x_0) = 0$, the
trivial group, since any two loops $f_0$ and $f_1$ based at $x_0$ are homotopic via the
linear homotopy of \cref{ex:linear-homotopies-convex}.
\end{example}

It is not so easy to show that a space has a nontrivial fundamental group since one
must somehow demonstrate the nonexistence of homotopies between certain loops.
We will tackle the simplest example shortly, computing the fundamental group of
the circle.

It is natural to ask about the dependence of $\pi_1(X,x_0)$ on the choice of basepoint.
Since $\pi_1(X,x_0)$ involves only the path-component of $X$ containing $x_0$, a relation
between $\pi_1(X,x_0)$ and $\pi_1(X,x_1)$ can be hoped for only if $x_0,x_1$ lie in the same
path-component. So let $h : I \to X$ be a path from $x_0$ to $x_1$, with inverse path
$\bar h(s) = h(1-s)$. We can then associate to each loop $f$ based at $x_1$ the loop
$h \cdot f \cdot \bar h$ based at $x_0$.



\begin{definition}[Change-of-basepoint map]
\label{def:change-of-basepoint-map}
\group{Fundamental Group}
\uses{def:path-composition,def:loop-and-pi1-set}
For a path $h$ from $x_0$ to $x_1$, the \textbf{change-of-basepoint} map
$\beta_h : \pi_1(X,x_1) \to \pi_1(X,x_0)$ is defined by $\beta_h[f] = [h \cdot f \cdot \bar h]$ (the two ways of
associating the triple product give homotopic results). This is well-defined since if
$f_t$ is a homotopy of loops based at $x_1$ then $h \cdot f_t \cdot \bar h$ is a homotopy of loops
based at $x_0$.
\end{definition}

\begin{proposition}[The change-of-basepoint map is an isomorphism]
\label{prop:change-of-basepoint-isomorphism}
\dcref{ch1:1.5}
\group{Fundamental Group}
\uses{def:change-of-basepoint-map}
The map $\beta_h : \pi_1(X,x_1) \to \pi_1(X,x_0)$ is an isomorphism.
\end{proposition}

\begin{proof}
First, $\beta_h$ is a homomorphism since $\beta_h[f\cdot g] = [h\cdot f\cdot g\cdot\bar h]
= [h\cdot f\cdot\bar h\cdot h\cdot g\cdot\bar h] = \beta_h[f]\beta_h[g]$. Further, $\beta_h$ is an isomorphism with
inverse $\beta_{\bar h}$ since $\beta_h\beta_{\bar h}[f] = \beta_h[\bar h\cdot f\cdot h] = [h\cdot\bar h\cdot f\cdot h\cdot\bar h] = [f]$, and
similarly $\beta_{\bar h}\beta_h[f]=[f]$.
\end{proof}

Thus if $X$ is path-connected, $\pi_1(X,x_0)$ is, up to isomorphism, independent of the
choice of basepoint, and the notation is often abbreviated $\pi_1(X)$ or $\pi_1 X$.

\begin{definition}[Simply-connected space]
\label{def:simply-connected}
\group{Fundamental Group}
\uses{def:loop-and-pi1-set}
A space is \textbf{simply-connected} if it is path-connected and has trivial
fundamental group.
\end{definition}

\begin{proposition}[Simply-connected iff unique homotopy class of connecting paths]
\label{prop:simply-connected-iff-unique-path-class}
\dcref{ch1:1.6}
\group{Fundamental Group}
\uses{def:simply-connected,def:homotopy-of-paths}
A space $X$ is simply-connected iff there is a unique homotopy class of paths
connecting any two points in $X$.
\end{proposition}

\begin{proof}
Path-connectedness is the existence of connecting paths, so only uniqueness needs
proof. Suppose $\pi_1(X) = 0$. If $f,g$ are paths from $x_0$ to $x_1$, then
$f \simeq f\cdot\bar g\cdot g = g$ since $\bar g\cdot g$ and $f\cdot\bar g$ (a loop at $x_0$) are each
homotopic to constant loops. Conversely, if there is only one homotopy class of
paths connecting $x_0$ to itself, then all loops at $x_0$ are homotopic to the constant
loop and $\pi_1(X,x_0) = 0$.
\end{proof}

\subsection*{The Fundamental Group of the Circle}

Our first real theorem will be the calculation $\pi_1(S^1) \cong \Z$.

\begin{theorem}[The fundamental group of the circle]
\label{thm:fundamental-group-of-circle}
\dcref{ch1:1.7}
\group{Fundamental Group}
\uses{def:loop-and-pi1-set,def:covering-space}
$\pi_1(S^1)$ is an infinite cyclic group generated by the homotopy class of the loop
$\omega(s) = (\cos 2\pi s, \sin 2\pi s)$ based at $(1,0)$.
\end{theorem}

Note that $[\omega]^n = [\omega_n]$ where $\omega_n(s) = (\cos 2\pi n s, \sin 2\pi n s)$ for $n \in \Z$. The
theorem is equivalent to: every loop in $S^1$ based at $(1,0)$ is homotopic to $\omega_n$ for a
unique $n \in \Z$. The idea is to compare paths in $S^1$ with paths in $\R$ via the map
$p : \R \to S^1$, $p(s) = (\cos 2\pi s, \sin 2\pi s)$, visualized as projecting the helix
$s \mapsto (\cos 2\pi s, \sin 2\pi s, s) \subset \R^3$ onto $\R^2$. Observe $\omega_n = p\tilde\omega_n$ where
$\tilde\omega_n(s) = ns$; the relation $\omega_n = p\tilde\omega_n$ is expressed by saying $\tilde\omega_n$ is a \textbf{lift} of
$\omega_n$.



\begin{definition}[Covering space]
\label{def:covering-space}
\group{Fundamental Group}
Given a space $X$, a \textbf{covering space} of $X$ consists of a space $\tilde X$ and a map
$p : \tilde X \to X$ such that for each point $x \in X$ there is an open neighborhood $U$ of
$x$ in $X$, called \textbf{evenly covered}, such that $p^{-1}(U)$ is a union of disjoint open sets
each mapped homeomorphically onto $U$ by $p$.
\end{definition}

For example, for $p : \R \to S^1$ above, any open arc in $S^1$ is evenly covered. To
prove \cref{thm:fundamental-group-of-circle} we need just two facts about covering spaces
$p : \tilde X \to X$:
\begin{enumerate}
\item[(a)] For each path $f : I \to X$ starting at $x_0 \in X$ and each $\tilde x_0 \in p^{-1}(x_0)$
there is a unique lift $\tilde f : I \to \tilde X$ starting at $\tilde x_0$.
\item[(b)] For each homotopy $f_t : I \to X$ of paths starting at $x_0$ and each
$\tilde x_0 \in p^{-1}(x_0)$ there is a unique lifted homotopy $\tilde f_t : I \to \tilde X$ of paths starting
at $\tilde x_0$.
\end{enumerate}

\begin{proof}[Proof of \cref{thm:fundamental-group-of-circle}]
Let $f : I \to S^1$ be a loop at $x_0 = (1,0)$. By (a) there is a lift $\tilde f$ starting at
$0$, ending at some integer $n$ since $p\tilde f(1) = f(1) = x_0$ and $p^{-1}(x_0) = \Z \subset \R$.
Another path in $\R$ from $0$ to $n$ is $\tilde\omega_n$, and $\tilde f \simeq \tilde\omega_n$ via the linear homotopy
$(1-t)\tilde f + t\tilde\omega_n$. Composing with $p$ gives $f \simeq \omega_n$, so $[f] = [\omega_n]$.

To show $n$ is uniquely determined by $[f]$: suppose $f \simeq \omega_n$ and $f \simeq \omega_m$, so
$\omega_m \simeq \omega_n$ via some $f_t$. By (b) this lifts to a homotopy $\tilde f_t$ of paths starting at
$0$; uniqueness in (a) gives $\tilde f_0 = \tilde\omega_m$, $\tilde f_1 = \tilde\omega_n$. Since $\tilde f_t$ is a homotopy of
paths, the endpoint $\tilde f_t(1)$ is independent of $t$, so $m = n$.

It remains to prove (a) and (b), both deduced from:
\begin{enumerate}
\item[(c)] Given $F : Y \times I \to X$ and a lift $\tilde F : Y \times \{0\} \to \tilde X$ of
$F|_{Y\times\{0\}}$, there is a unique lift $\tilde F : Y \times I \to \tilde X$ extending it.
\end{enumerate}
Statement (a) is the case $Y$ a point; (b) follows by applying (c) with $Y = I$ to
$F(s,t) = f_t(s)$, using (a) to lift $F|_{I \times \{0\}}$, and observing that the restrictions of
the resulting $\tilde F$ to $\{0\}\times I$ and $\{1\}\times I$ must be constant (they lift constant
paths), so $\tilde f_t(s) = \tilde F(s,t)$ is a homotopy of paths lifting $f_t$.

For (c): fixing $y_0 \in Y$, continuity of $F$ and compactness of $\{y_0\}\times I$ give a
neighborhood $N$ of $y_0$ and a partition $0=t_0<\cdots<t_m=1$ with each
$F(N\times[t_i,t_{i+1}])$ inside an evenly covered $U_i$. Inductively lift $\tilde F$ over
$N\times[t_i,t_{i+1}]$ by composing $F$ with $p^{-1}: U_i \to \tilde U_i$ for the sheet $\tilde U_i$
containing $\tilde F(y_0,t_i)$, shrinking $N$ as needed so $\tilde F(N\times\{t_i\}) \subset \tilde U_i$. This
gives a lift on some $N \times I$.

Uniqueness (case $Y$ a point, so $\tilde f, \tilde f'$ lifting $f$ with $\tilde f(0)=\tilde f'(0)$): with the
same partition, if $\tilde f = \tilde f'$ on $[0,t_i]$ then $\tilde f([t_i,t_{i+1}])$ and $\tilde f'([t_i,t_{i+1}])$,
being connected and agreeing at $t_i$, lie in the same sheet $\tilde U_i$, on which $p$ is
injective and $p\tilde f = p\tilde f' = f$, so $\tilde f = \tilde f'$ there too.

Finally, the local lifts on sets $N \times I$ agree on overlaps by the uniqueness just
proved (restricted to each segment $\{y\}\times I$), so they patch to a well-defined,
continuous, unique lift $\tilde F$ on all of $Y \times I$.
\end{proof}

Now we turn to some applications of the calculation of $\pi_1(S^1)$.

\begin{theorem}[Fundamental Theorem of Algebra]
\label{thm:fundamental-theorem-of-algebra}
\dcref{ch1:1.8}
\group{Fundamental Group}
\uses{thm:fundamental-group-of-circle}
Every nonconstant polynomial with coefficients in $\C$ has a root in $\C$.
\end{theorem}

\begin{proof}
We may assume $p(z) = z^n + a_1z^{n-1} + \cdots + a_n$. If $p(z)$ has no roots in $\C$,
then for each $r \ge 0$,
\[
f_r(s) = \frac{p(re^{2\pi i s})/p(r)}{|p(re^{2\pi i s})/p(r)|}
\]
defines a loop in $S^1 \subset \C$ based at $1$, and as $r$ varies, $f_r$ is a homotopy of
loops based at $1$. Since $f_0$ is trivial, $[f_r] = 0$ in $\pi_1(S^1)$ for all $r$. Fix $r$ larger
than $1$ and than $|a_1|+\cdots+|a_n|$. Then for $|z|=r$,
\[
|z^n| > (|a_1|+\cdots+|a_n|)|z|^{n-1} \ge |a_1z^{n-1}+\cdots+a_n|,
\]
so $p_t(z) = z^n + t(a_1z^{n-1}+\cdots+a_n)$ has no roots on $|z|=r$ for $0 \le t \le 1$.
Replacing $p$ by $p_t$ in $f_r$ and letting $t$ go from $1$ to $0$ gives a homotopy from
$f_r$ to $\omega_n(s) = e^{2\pi i n s}$. By \cref{thm:fundamental-group-of-circle}, $\omega_n$ represents $n$ times a
generator of $\pi_1(S^1)$; since $[\omega_n] = [f_r] = 0$, we get $n=0$. Thus the only
polynomials without roots in $\C$ are constants.
\end{proof}

\begin{theorem}[Brouwer fixed point theorem, dimension 2]
\label{thm:brouwer-fixed-point-dim2}
\dcref{ch1:1.9}
\group{Fundamental Group}
\uses{thm:fundamental-group-of-circle,def:retraction}
Every continuous map $h : D^2 \to D^2$ has a fixed point.
\end{theorem}

Here $D^n$ denotes the closed unit disk in $\R^n$, with boundary the unit sphere
$S^{n-1}$.

\begin{proof}
Suppose $h(x) \ne x$ for all $x \in D^2$. Define $r : D^2 \to S^1$ by letting $r(x)$ be the
point of $S^1$ where the ray from $h(x)$ through $x$ leaves $D^2$; $r$ is continuous and
$r(x) = x$ for $x \in S^1$, so $r$ is a retraction of $D^2$ onto $S^1$.



Let $f_0$ be any loop in $S^1$; in $D^2$, $f_0$ is homotopic to a constant loop via
$f_t(s) = (1-t)f_0(s)+tx_0$. Since $r$ is the identity on $S^1$, $r\circ f_t$ is a homotopy in
$S^1$ from $f_0$ to a constant loop, contradicting $\pi_1(S^1) \ne 0$.
\end{proof}

This theorem was first proved by Brouwer around 1910. Brouwer in fact proved the
corresponding result for $D^n$ (obtained in \cref{cor:non-retraction-disk-boundary} using homology), using the notion of degree, which is defined in Section 2.2. A more
constructive proof is due to Sperner in 1928, explained in \cite{AignerZiegler1999}.

\begin{theorem}[Borsuk--Ulam theorem, dimension 2]
\label{thm:borsuk-ulam-dim2}
\dcref{ch1:1.10}
\group{Fundamental Group}
\uses{thm:fundamental-group-of-circle}
For every continuous map $f : S^2 \to \R^2$ there exists a pair of antipodal points $x$
and $-x$ in $S^2$ with $f(x) = f(-x)$.
\end{theorem}

The theorem holds more generally for maps $S^n \to \R^n$ (\cref{cor:borsuk-ulam}). The case $n=1$ is easy since $f(x)-f(-x)$ changes sign as $x$ goes
halfway around the circle.

\begin{proof}
If false for some $f : S^2 \to \R^2$, define $g : S^2 \to S^1$ by
$g(x) = (f(x)-f(-x))/|f(x)-f(-x)|$. Let $\eta(s) = (\cos 2\pi s, \sin 2\pi s, 0)$ be the
equatorial loop and $h = g\eta$. Since $g(-x) = -g(x)$, we get $h(s+1/2) = -h(s)$ for
$s \in [0,1/2]$. Lift $h$ to $\tilde h : I \to \R$; the relation implies
$\tilde h(s+1/2) = \tilde h(s) + q/2$ for some odd integer $q$, and $q$ is independent of $s$
since it depends continuously on $s$ but takes integer values. In particular
$\tilde h(1) = \tilde h(0) + q$, so $h$ represents $q$ times a generator of $\pi_1(S^1)$; since $q$ is
odd, $h$ is not nullhomotopic. But $\eta$ is nullhomotopic in $S^2$, so $h = g\eta$ is
nullhomotopic in $S^1$ -- a contradiction.
\end{proof}

\begin{corollary}[Covering $S^2$ by three closed sets]
\label{cor:sphere-three-closed-sets-antipodal}
\dcref{ch1:1.11}
\group{Fundamental Group}
\uses{thm:borsuk-ulam-dim2}
Whenever $S^2$ is expressed as the union of three closed sets $A_1,A_2,A_3$, at least
one of these sets contains a pair of antipodal points $\{x,-x\}$.
\end{corollary}

\begin{proof}
Let $d_i(x) = \inf_{y\in A_i}|x-y|$, a continuous function. Applying
\cref{thm:borsuk-ulam-dim2} to $x \mapsto (d_1(x),d_2(x))$ gives antipodal $x,-x$ with
$d_1(x)=d_1(-x)$ and $d_2(x)=d_2(-x)$. If either distance is zero, $x$ and $-x$ lie in the
same closed set $A_1$ or $A_2$; otherwise both distances are positive, so $x,-x \notin A_1
\cup A_2$ and hence both lie in $A_3$.
\end{proof}

To see that `three' is best possible, project the four faces of a tetrahedron
circumscribing $S^2$ radially onto the sphere, giving a cover by four closed sets none
containing an antipodal pair. The same argument shows $S^n$ cannot be covered by
$n+1$ closed sets without antipodal pairs (assuming the higher-dimensional
Borsuk--Ulam theorem), though it can be covered by $n+2$ such sets.

\begin{proposition}[$\pi_1$ of a product]
\label{prop:pi1-of-product-space}
\dcref{ch1:1.12}
\group{Fundamental Group}
\uses{def:loop-and-pi1-set}
$\pi_1(X\times Y)$ is isomorphic to $\pi_1(X)\times\pi_1(Y)$ if $X$ and $Y$ are
path-connected.
\end{proposition}

\begin{proof}
A map $f : Z \to X\times Y$ is continuous iff its two coordinate maps are, so a loop
in $X\times Y$ at $(x_0,y_0)$ is equivalent to a pair of loops in $X$ and $Y$, and likewise
for homotopies. This gives a bijection
$\pi_1(X\times Y,(x_0,y_0)) \approx \pi_1(X,x_0)\times\pi_1(Y,y_0)$, $[f]\mapsto([g],[h])$, obviously a
group homomorphism and hence an isomorphism.
\end{proof}

\begin{example}[The torus]
\label{ex:torus-fundamental-group}
\dcref{ch1:1.13}
\group{Fundamental Group}
\uses{prop:pi1-of-product-space,thm:fundamental-group-of-circle}
By \cref{prop:pi1-of-product-space}, $\pi_1(S^1\times S^1) \approx \Z\times\Z$. A pair
$(p,q) \in \Z\times\Z$ corresponds to a loop $\omega_{p,q}(s) = (\omega_p(s),\omega_q(s))$ winding $p$
times around one $S^1$ factor and $q$ times around the other. This loop can be
knotted, as for $p=3,q=2$ -- the resulting torus knots are studied later
(\cref{ex:torus-knot-complement-fundamental-group}).



More generally, the $n$-dimensional torus (product of $n$ circles) has fundamental
group isomorphic to $\Z^n$, by induction on $n$.
\end{example}

\subsection*{Induced Homomorphisms}

\begin{definition}[Induced homomorphism]
\label{def:induced-homomorphism}
\group{Fundamental Group}
\uses{def:loop-and-pi1-set}
Suppose $\varphi : X \to Y$ takes basepoint $x_0 \in X$ to basepoint $y_0 \in Y$, written
$\varphi : (X,x_0) \to (Y,y_0)$. Then $\varphi$ \textbf{induces} a homomorphism
$\varphi_* : \pi_1(X,x_0) \to \pi_1(Y,y_0)$ by $\varphi_*[f] = [\varphi f]$.
\end{definition}

This is well-defined since a homotopy $f_t$ of loops at $x_0$ yields a composed
homotopy $\varphi f_t$ of loops at $y_0$; it is a homomorphism since
$\varphi(f\cdot g) = (\varphi f)\cdot(\varphi g)$. Two basic properties: $(\varphi\psi)_* = \varphi_*\psi_*$ for a
composition $(X,x_0)\xrightarrow{\psi}(Y,y_0)\xrightarrow{\varphi}(Z,z_0)$, and $\one_* = \one$. These make
the fundamental group a functor (the formal definition of `functor' is defined in
\cref{sec:formal-viewpoint}).

As an application, if $\varphi$ is a homeomorphism with inverse $\psi$ then $\varphi_*$ is an
isomorphism with inverse $\psi_*$, since $\varphi_*\psi_* = (\varphi\psi)_* = \one_* = \one$ and likewise
$\psi_*\varphi_* = \one$.

\begin{proposition}[$\pi_1(S^n) = 0$ for $n \ge 2$]
\label{prop:pi1-sphere-trivial-n-geq-2}
\dcref{ch1:1.14}
\group{Fundamental Group}
\uses{def:induced-homomorphism,lem:loop-decomposition-open-cover,ex:convex-set-trivial-pi1}
$\pi_1(S^n) = 0$ if $n \ge 2$.
\end{proposition}

The main step will be a general fact also used in \cref{sec:van-kampen-theorem}.

\begin{lemma}[Loops decompose over a path-connected open cover]
\label{lem:loop-decomposition-open-cover}
\dcref{ch1:1.15}
\group{Fundamental Group}
\uses{def:loop-and-pi1-set}
If a space $X$ is the union of path-connected open sets $A_\alpha$ each containing the
basepoint $x_0 \in X$, and each intersection $A_\alpha \cap A_\beta$ is path-connected, then every
loop in $X$ at $x_0$ is homotopic to a product of loops each contained in a single
$A_\alpha$.
\end{lemma}

\begin{proof}
Given a loop $f : I \to X$ at $x_0$, continuity and compactness of $I$ give a partition
$0=s_0<\cdots<s_m=1$ such that each $f([s_{i-1},s_i])$ lies in a single $A_i$. Let $f_i$ be
$f$ restricted to $[s_{i-1},s_i]$, so $f = f_1\cdot\cdots\cdot f_m$. Since $A_i\cap A_{i+1}$ is
path-connected, choose a path $g_i$ in $A_i\cap A_{i+1}$ from $x_0$ to $f(s_i)$. Then
\[
f \simeq (f_1\cdot\bar g_1)\cdot(g_1\cdot f_2\cdot\bar g_2)\cdot(g_2\cdot f_3\cdot\bar g_3)\cdots(g_{m-1}\cdot f_m),
\]
a composition of loops each lying in a single $A_i$.
\end{proof}



\begin{proof}[Proof of \cref{prop:pi1-sphere-trivial-n-geq-2}]
Write $S^n = A_1 \cup A_2$ with $A_1,A_2$ each homeomorphic to $\R^n$ (complements of
two antipodal points) and $A_1\cap A_2$ homeomorphic to $S^{n-1}\times\R$. Choose the
basepoint in $A_1\cap A_2$; for $n\ge 2$ this is path-connected, so
\cref{lem:loop-decomposition-open-cover} says every loop in $S^n$ is homotopic to a product of
loops in $A_1$ or $A_2$, both of which have trivial $\pi_1$ since $A_1,A_2 \approx \R^n$. Hence
every loop in $S^n$ is nullhomotopic.
\end{proof}

\begin{corollary}[$\R^2$ is not homeomorphic to $\R^n$ for $n \ne 2$]
\label{cor:r2-not-homeomorphic-rn}
\dcref{ch1:1.16}
\group{Fundamental Group}
\uses{prop:pi1-of-product-space,prop:pi1-sphere-trivial-n-geq-2,thm:fundamental-group-of-circle}
$\R^2$ is not homeomorphic to $\R^n$ for $n \ne 2$.
\end{corollary}

\begin{proof}
Suppose $f : \R^2 \to \R^n$ is a homeomorphism. For $n=1$, $\R^2-\{0\}$ is
path-connected but $\R^n-\{f(0)\}$ is not. For $n>2$, use fundamental groups: for
$x\in\R^n$, $\R^n-\{x\} \approx S^{n-1}\times\R$, so by \cref{prop:pi1-of-product-space},
$\pi_1(\R^n-\{x\}) \approx \pi_1(S^{n-1})$, which is $\Z$ for $n=2$ (\cref{thm:fundamental-group-of-circle})
and trivial for $n>2$ (\cref{prop:pi1-sphere-trivial-n-geq-2}).
\end{proof}

The more general fact that $\R^m \not\approx \R^n$ for $m \ne n$ is proved the same way
using higher homology groups (\cref{thm:invariance-of-dimension}).

\begin{proposition}[Retracts and deformation retracts induce injections/isomorphisms on $\pi_1$]
\label{prop:retract-induces-injective-pi1}
\dcref{ch1:1.17}
\group{Fundamental Group}
\uses{def:induced-homomorphism,def:retraction,def:deformation-retraction}
If a space $X$ retracts onto a subspace $A$, then the homomorphism
$i_* : \pi_1(A,x_0) \to \pi_1(X,x_0)$ induced by the inclusion $i : A \hookrightarrow X$ is injective.
If $A$ is a deformation retract of $X$, then $i_*$ is an isomorphism.
\end{proposition}

\begin{proof}
If $r : X \to A$ is a retraction, $r\circ i = \one$, so $r_*i_* = \one$, giving $i_*$ injective. If
$r_t$ is a deformation retraction of $X$ onto $A$, then for a loop $f$ at $x_0 \in A$,
$r_1 \circ f$ gives a homotopy of $f$ to a loop in $A$, so $i_*$ is also surjective.
\end{proof}

This gives another way of seeing that $S^1$ is not a retract of $D^2$, since the
inclusion-induced map $\pi_1(S^1)\to\pi_1(D^2)$ is $\Z\to 0$, which cannot be injective.

\begin{definition}[Homotopy of maps of pairs; based homotopy equivalence]
\label{def:based-homotopy-equivalence}
\group{Fundamental Group}
\uses{def:homotopy}
Recall from Chapter 0 (\cref{def:homotopy}) the notion of a homotopy $\varphi_t : X \to Y$.
If $\varphi_t$ takes a subspace $A\subset X$ into $B\subset Y$ for all $t$, we speak of a homotopy
of maps of pairs $\varphi_t : (X,A) \to (Y,B)$; if additionally $\varphi_t(x_0)=y_0$ for all $t$ it is a
\textbf{basepoint-preserving homotopy} $\varphi_t : (X,x_0)\to(Y,y_0)$. One says
$(X,x_0) \simeq (Y,y_0)$ if there are $\varphi : (X,x_0)\to(Y,y_0)$ and $\psi : (Y,y_0)\to(X,x_0)$ with
$\varphi\psi \simeq \one$ and $\psi\varphi \simeq \one$ through basepoint-fixing homotopies.
\end{definition}

If $\varphi_t : (X,x_0)\to(Y,y_0)$ is a basepoint-preserving homotopy, then
$\varphi_{0*} = \varphi_{1*}$, since $\varphi_{0*}[f] = [\varphi_0 f] = [\varphi_1 f] = \varphi_{1*}[f]$ via the homotopy
$\varphi_t f$. Consequently if $(X,x_0)\simeq(Y,y_0)$ then $\varphi_*\psi_* = (\varphi\psi)_* = \one$ and
$\psi_*\varphi_* = \one$, so $\varphi_*,\psi_*$ are inverse isomorphisms $\pi_1(X,x_0)\cong\pi_1(Y,y_0)$ -- in
particular this reproves \cref{prop:retract-induces-injective-pi1} for deformation retracts,
since $X$ deformation retracting onto $A$ gives $(X,x_0)\simeq(A,x_0)$.

\begin{proposition}[Homotopy equivalences induce isomorphisms on $\pi_1$, basepoints unconstrained]
\label{prop:homotopy-equivalence-induces-pi1-iso}
\dcref{ch1:1.18}
\group{Fundamental Group}
\uses{def:induced-homomorphism,lem:basepoint-change-homotopy,def:homotopy-equivalence}
If $\varphi : X \to Y$ is a homotopy equivalence, then
$\varphi_* : \pi_1(X,x_0) \to \pi_1(Y,\varphi(x_0))$ is an isomorphism for all $x_0 \in X$.
\end{proposition}

\begin{lemma}[Basepoint tracking under an unbased homotopy]
\label{lem:basepoint-change-homotopy}
\dcref{ch1:1.19}
\group{Fundamental Group}
\uses{def:change-of-basepoint-map,def:induced-homomorphism}
If $\varphi_t : X \to Y$ is a homotopy and $h$ is the path $\varphi_t(x_0)$ traced by a basepoint
$x_0 \in X$, then $\varphi_{0*} = \beta_h \varphi_{1*} : \pi_1(X,x_0) \to \pi_1(Y,\varphi_0(x_0))$.
\end{lemma}



\begin{proof}
Let $h_t(s) = h(ts)$, the restriction of $h$ to $[0,t]$ reparametrized to $[0,1]$. For a
loop $f$ in $X$ at $x_0$, the product $h_t\cdot(\varphi_t f)\cdot\bar h_t$ gives a homotopy of loops at
$\varphi_0(x_0)$; restricting to $t=0$ and $t=1$ gives $\varphi_{0*}[f] = \beta_h(\varphi_{1*}[f])$.
\end{proof}



\begin{proof}[Proof of \cref{prop:homotopy-equivalence-induces-pi1-iso}]
Let $\psi : Y \to X$ be a homotopy inverse for $\varphi$, so $\varphi\psi\simeq\one$, $\psi\varphi\simeq\one$.
Consider
\[
\pi_1(X,x_0) \xrightarrow{\varphi_*} \pi_1(Y,\varphi(x_0)) \xrightarrow{\psi_*} \pi_1(X,\psi\varphi(x_0))
\xrightarrow{\varphi_*} \pi_1(Y,\varphi\psi\varphi(x_0)).
\]
The composition of the first two maps is $\psi_*\varphi_* = \beta_h$ for some $h$, by
\cref{lem:basepoint-change-homotopy}, hence an isomorphism; in particular $\varphi_*$ is
injective. The same reasoning applied to the second and third maps shows $\psi_*$ is
injective. So the first two maps are injective with an isomorphic composition, hence
the first map $\varphi_*$ is also surjective, so an isomorphism.
\end{proof}

\subsection*{Exercises}
% Exercises are transcribed for completeness but are not given blueprint labels or
% \uses edges.

\begin{enumerate}
\item Show that composition of paths satisfies the following cancellation property:
if $f_0\cdot g_0 = f_1\cdot g_1$ and $g_0=g_1$ then $f_0=f_1$.

\item Show that the change-of-basepoint homomorphism $\beta_h$ depends only on the
homotopy class of $h$.

\item For a path-connected space $X$, show that $\pi_1(X)$ is abelian iff all
basepoint-change homomorphisms $\beta_h$ depend only on the endpoints of $h$.

\item A subspace $X\subset\R^n$ is \textit{star-shaped} if there is $x_0\in X$ such that for
each $x\in X$ the segment from $x_0$ to $x$ lies in $X$. Show that if $X$ is locally
star-shaped, every path in $X$ is homotopic in $X$ to a piecewise linear path. Show
this applies when $X$ is open or a union of finitely many closed convex sets.

\item Show that for a space $X$, the following are equivalent: (a) every map
$S^1\to X$ is homotopic to a constant map; (b) every map $S^1\to X$ extends to a map
$D^2\to X$; (c) $\pi_1(X,x_0)=0$ for all $x_0\in X$. Deduce that $X$ is simply-connected
iff all maps $S^1\to X$ are homotopic (without regard to basepoints).

\item Let $[S^1,X]$ be the set of homotopy classes of maps $S^1\to X$ with no
basepoint condition, and $\Phi:\pi_1(X,x_0)\to[S^1,X]$ the natural map. Show $\Phi$ is
onto if $X$ is path-connected, and $\Phi([f])=\Phi([g])$ iff $[f],[g]$ are conjugate in
$\pi_1(X,x_0)$; deduce a bijection between $[S^1,X]$ and conjugacy classes in $\pi_1(X)$.

\item Define $f:S^1\times I\to S^1\times I$ by $f(\theta,s)=(\theta+2\pi s,s)$. Show $f$ is
homotopic to the identity by a homotopy stationary on one boundary circle, but not
by any homotopy stationary on both.

\item Does the Borsuk--Ulam theorem hold for the torus, i.e. must every map
$f:S^1\times S^1\to\R^2$ satisfy $f(x,y)=f(-x,-y)$ for some $(x,y)$?

\item Let $A_1,A_2,A_3$ be compact sets in $\R^3$. Use Borsuk--Ulam to show there is a
single plane $P\subset\R^3$ simultaneously bisecting the measure of each $A_i$.

\item From $\pi_1(X\times Y,(x_0,y_0)) \simeq \pi_1(X,x_0)\times\pi_1(Y,y_0)$, loops in
$X\times\{y_0\}$ and $\{x_0\}\times Y$ represent commuting elements. Construct an explicit
homotopy demonstrating this.

\item If $X_0$ is the path-component of $X$ containing $x_0$, show the inclusion
$X_0\hookrightarrow X$ induces an isomorphism $\pi_1(X_0,x_0)\to\pi_1(X,x_0)$.

\item Show that every homomorphism $\pi_1(S^1)\to\pi_1(S^1)$ is induced by some map
$\varphi:S^1\to S^1$.

\item Given $X$ and a path-connected subspace $A$ containing $x_0$, show
$\pi_1(A,x_0)\to\pi_1(X,x_0)$ is surjective iff every path in $X$ with endpoints in $A$ is
homotopic to a path in $A$.

\item Show that the isomorphism $\pi_1(X\times Y)\simeq\pi_1(X)\times\pi_1(Y)$ of
\cref{prop:pi1-of-product-space} is $[f]\mapsto(p_{1*}[f],p_{2*}[f])$ for the two projections.

\item Given $f:X\to Y$ and a path $h:I\to X$ from $x_0$ to $x_1$, show
$f_*\beta_h = \beta_{fh}f_*$.

\item Show there are no retractions $r:X\to A$ in the following cases: (a)
$X=\R^3$, $A$ any subspace homeomorphic to $S^1$; (b) $X=S^1\times D^2$, $A$ its
boundary torus $S^1\times S^1$; (c) $X=S^1\times D^2$, $A$ a certain circle in $X$; (d)
$X=D^2\vee D^2$, $A=S^1\vee S^1$ its boundary; (e) $X$ a disk with two boundary
points identified, $A=S^1\vee S^1$ its boundary; (f) $X$ the M\"obius band, $A$ its
boundary circle.

\item Construct infinitely many nonhomotopic retractions $S^1\vee S^1\to S^1$.

\item Using \cref{lem:loop-decomposition-open-cover}, show that if $X$ is obtained from a
path-connected subspace $A$ by attaching a cell $e^n$ with $n\ge 2$, then
$A\hookrightarrow X$ induces a surjection on $\pi_1$. Apply this to show: (a) $S^1\vee S^2$
has fundamental group $\Z$; (b) for a path-connected CW complex $X$, the inclusion
of the 1-skeleton $X^1\hookrightarrow X$ induces a surjection $\pi_1(X^1)\to\pi_1(X)$.

\item Show that if $X$ is a path-connected 1-dimensional CW complex with
basepoint a 0-cell, every loop in $X$ is homotopic to a loop consisting of finitely
many edges traversed monotonically. [This gives an elementary proof that
$\pi_1(S^1)$ is cyclic, generated by the standard loop winding once around $S^1$.]
\end{enumerate}

\section{1.2 Van Kampen's Theorem}
\label{sec:van-kampen-theorem}

The van Kampen theorem gives a method for computing the fundamental groups
of spaces that can be decomposed into simpler spaces whose fundamental groups are
already known. We shall see for example that for every group $G$ there is a space
$X_G$ whose fundamental group is isomorphic to $G$.

Consider the space $X$ formed by two circles $A$ and $B$ intersecting in a single
basepoint $x_0$. We know $\pi_1(A) \cong \Z$, generated by a loop $a$ once around $A$,
and likewise $\pi_1(B) \cong \Z$ generated by $b$. Each product of powers of $a$ and $b$,
such as $a^5b^2a^{-3}ba^2$, gives an element of $\pi_1(X)$. The set of all such alternating
words forms a group denoted $\Z * \Z$, with multiplication by juxtaposition and
reduction (e.g. $(b^4a^3b^2a^{-3})(a^3b^{-1}ab^3) = b^4a^3b^2b^{-1}ab^3 = b^4a^3bab^3$). It would be very nice if
$\pi_1(X) \cong \Z * \Z$ -- van Kampen's theorem will imply that this is indeed the case,
and more generally that a union of $k$ circles touching at a point has fundamental
group $\Z * \cdots * \Z$ ($k$ times).



The group $\Z * \Z$ is an example of the \textbf{free product} of groups, which we
describe next before stating van Kampen's theorem itself.

\subsection*{Free Products of Groups}

Given a collection of groups $G_\alpha$, the product $\prod_\alpha G_\alpha$ or direct sum
$\bigoplus_\alpha G_\alpha$ contain all the $G_\alpha$ as subgroups but force elements of different
$G_\alpha$'s to commute -- unnatural for nonabelian groups. The free product
$*_\alpha G_\alpha$ is the nonabelian analog of $\bigoplus_\alpha G_\alpha$.

\begin{definition}[Free product of groups]
\label{def:free-product-of-groups}
\group{Fundamental Group}
As a set, the \textbf{free product} $*_\alpha G_\alpha$ consists of all \textbf{reduced} words
$g_1g_2\cdots g_m$ of finite length $m \ge 0$, where each $g_i$ belongs to some $G_{\alpha_i}$,
$g_i \ne e$, and adjacent letters belong to different groups ($\alpha_i \ne \alpha_{i+1}$). The empty
word is the identity. The group operation is juxtaposition followed by reduction:
combine adjacent letters from the same $G_\alpha$ via that group's multiplication, and
cancel any resulting identity letters, repeating until reduced.
\end{definition}

Associativity of this operation is proved indirectly: for $g \in G_\alpha$, left
multiplication $L_g : W \to W$ (where $W$ is the set of reduced words) satisfies
$L_{gg'} = L_gL_{g'}$, giving an injective homomorphism $L : W \to P(W)$ into the
permutation group of $W$ by $L(g_1\cdots g_m) = L_{g_1}\cdots L_{g_m}$; since composition in
$P(W)$ is associative and the product on $W$ corresponds to composition under $L$,
the product on $W$ is associative.

\begin{definition}[Free group, basis, rank]
\label{def:free-group-basis-rank}
\group{Fundamental Group}
\uses{def:free-product-of-groups}
A \textbf{free group} is the free product of any number of copies of $\Z$, finite or
infinite; the generators, one for each $\Z$ factor, are called a \textbf{basis}, and the
number of basis elements is the \textbf{rank}. The abelianization of a free group is free
abelian with the same basis, so the rank is well-defined independent of the choice
of basis.
\end{definition}

An example of a free product that is not a free group is $\Z_2 * \Z_2$, consisting of the
alternating words in $a,b$ with $a^2=b^2=e$. The homomorphism
$\varphi : \Z_2*\Z_2 \to \Z_2$ sending a word to its length mod 2 has kernel the infinite
cyclic subgroup generated by $ab$ (since $ba=(ab)^{-1}$), and $\Z_2*\Z_2$ is the semidirect
product of this $\Z$ and the $\Z_2$ generated by $a$, with $a(ab)a^{-1} = (ab)^{-1}$ -- the
\textbf{infinite dihedral group}.

For a general free product $*_\alpha G_\alpha$, each $G_\alpha$ embeds as a subgroup (empty word
plus nonidentity one-letter words), these subgroups otherwise disjoint and sharing
only the identity. Associativity implies every product $g_1\cdots g_m$ of elements of the
$G_\alpha$'s has a unique reduced form, independent of the order of reduction operations.

\begin{proposition}[Universal property of the free product]
\label{prop:free-product-universal-property}
\group{Fundamental Group}
\uses{def:free-product-of-groups}
Any collection of homomorphisms $\varphi_\alpha : G_\alpha \to H$ extends uniquely to a
homomorphism $\varphi : *_\alpha G_\alpha \to H$, by $\varphi(g_1\cdots g_n) = \varphi_{\alpha_1}(g_1)\cdots\varphi_{\alpha_n}(g_n)$.
\end{proposition}

For example, for $G*H$ the inclusions $G,H \hookrightarrow G\times H$ induce a surjective
homomorphism $G*H \to G\times H$, via \cref{prop:free-product-universal-property}.

\subsection*{The van Kampen Theorem}

Suppose $X$ is the union of path-connected open sets $A_\alpha$ each containing the
basepoint $x_0$. The inclusion-induced maps $j_\alpha : \pi_1(A_\alpha) \to \pi_1(X)$ extend, by
\cref{prop:free-product-universal-property}, to a homomorphism
$\Phi : *_\alpha \pi_1(A_\alpha) \to \pi_1(X)$. Writing $i_{\alpha\beta} : \pi_1(A_\alpha\cap A_\beta) \to \pi_1(A_\alpha)$ for the
inclusion-induced map, $j_\alpha i_{\alpha\beta} = j_\beta i_{\beta\alpha}$ (both induced by $A_\alpha\cap A_\beta \hookrightarrow X$), so
$\ker\Phi$ contains all $i_{\alpha\beta}(\omega)i_{\beta\alpha}(\omega)^{-1}$ for $\omega \in \pi_1(A_\alpha\cap A_\beta)$.

\begin{theorem}[Van Kampen's theorem]
\label{thm:van-kampen-theorem}
\dcref{ch1:1.20}
\group{Fundamental Group}
\uses{prop:free-product-universal-property,lem:loop-decomposition-open-cover}
If $X$ is the union of path-connected open sets $A_\alpha$ each containing the
basepoint $x_0$, and each $A_\alpha \cap A_\beta$ is path-connected, then
$\Phi : *_\alpha \pi_1(A_\alpha) \to \pi_1(X)$ is surjective. If in addition each
$A_\alpha \cap A_\beta \cap A_\gamma$ is path-connected, then $\ker\Phi$ is the normal subgroup $N$
generated by all $i_{\alpha\beta}(\omega)i_{\beta\alpha}(\omega)^{-1}$, and $\Phi$ induces an isomorphism
$\pi_1(X) \cong *_\alpha \pi_1(A_\alpha) / N$.
\end{theorem}

\begin{example}[Wedge sums]
\label{ex:wedge-sum-fundamental-group}
\dcref{ch1:1.21}
\group{Fundamental Group}
\uses{thm:van-kampen-theorem,def:wedge-sum}
If each basepoint $x_\alpha \in X_\alpha$ (\cref{def:wedge-sum}) is a deformation retract of an
open neighborhood $U_\alpha$, then $X_\alpha$ is a deformation retract of
$A_\alpha = X_\alpha \vee \bigvee_{\beta\ne\alpha} U_\beta$ inside $\bigvee_\alpha X_\alpha$, and any intersection of two or
more distinct $A_\alpha$'s is $\bigvee_\alpha U_\alpha$, contractible. \Cref{thm:van-kampen-theorem} then
gives that $\Phi : *_\alpha \pi_1(X_\alpha) \to \pi_1(\bigvee_\alpha X_\alpha)$ is an isomorphism when each $X_\alpha$
is path-connected. In particular $\pi_1(\bigvee_\alpha S^1_\alpha)$ is free on one generator per circle,
so $\pi_1(S^1\vee S^1) \cong \Z*\Z$. More generally the fundamental group of any
connected graph is free (\cref{prop:fundamental-group-of-graph-is-free}).
\end{example}

\begin{example}[A cube graph has free fundamental group on five generators]
\label{ex:cube-graph-fundamental-group-free}
\dcref{ch1:1.22}
\group{Fundamental Group}
\uses{thm:van-kampen-theorem}
Let $X$ be the graph of the twelve edges of a cube, and $T \subset X$ a maximal
tree (a contractible spanning subgraph) using seven of the edges. For each of the
five edges $e_\alpha$ of $X - T$, choose an open neighborhood $A_\alpha$ of $T\cup e_\alpha$ that
deformation retracts onto it; any intersection of two or more $A_\alpha$'s deformation
retracts onto $T$, hence is contractible, and each $A_\alpha \simeq S^1$. By
\cref{thm:van-kampen-theorem}, $\pi_1(X) \cong *_\alpha \pi_1(A_\alpha)$ is free on five generators, one
per edge outside $T$: for each such $e_\alpha$, a loop that goes through $T$ to one end of
$e_\alpha$, across $e_\alpha$, and back through $T$.


\end{example}

Van Kampen's theorem is often applied with just two sets $A_\alpha, A_\beta$, where the
triple-intersection hypothesis is vacuous, giving
$\pi_1(X) \cong (\pi_1(A_\alpha)*\pi_1(A_\beta))/N$. Path-connectedness of $A_\alpha\cap A_\beta$ is genuinely
needed: for $S^1$ decomposed into two open arcs, $\Phi$ is not surjective. Triple-
intersection path-connectedness is also needed: for $X$ the suspension of three
points $a,b,c$ with $A_\alpha,A_\beta,A_\gamma$ the complements of each point, the pairwise cover
$\{A_\alpha,A_\beta\}$ gives $\pi_1(X) \cong \Z*\Z$ (since $A_\alpha\cap A_\beta$ is contractible), but naively
applying the theorem to $\{A_\alpha,A_\beta,A_\gamma\}$ (pairwise but not triply intersection
path-connected) would wrongly suggest $\pi_1(X)\cong\Z*\Z*\Z$, which has a different
abelianization.



\begin{proof}[Proof of \cref{thm:van-kampen-theorem}]
Surjectivity of $\Phi$ is exactly \cref{lem:loop-decomposition-open-cover}. For the
harder half, call a \textbf{factorization} of $[f]\in\pi_1(X)$ a word $[f_1]\cdots[f_k]$ where
each $f_i$ is a loop in some $A_\alpha$ at $x_0$ and $f \simeq f_1\cdots f_k$ in $X$; surjectivity of
$\Phi$ says every $[f]$ has a factorization. Call two factorizations of $[f]$
\textbf{equivalent} if related by a sequence of: (i) combining adjacent terms in the same
$\pi_1(A_\alpha)$ into one term, or (ii) reinterpreting a term $[f_i]\in\pi_1(A_\alpha)$ as lying in
$\pi_1(A_\beta)$ when $f_i$ is a loop in $A_\alpha\cap A_\beta$. Move (i) does not change the
element of $*_\alpha\pi_1(A_\alpha)$ defined by the word; move (ii) does not change its image in
$Q = *_\alpha\pi_1(A_\alpha)/N$. So equivalent factorizations define the same element of $Q$.

It remains to show any two factorizations of $[f]$ are equivalent, which shows
$Q \to \pi_1(X)$ is injective and hence $\ker\Phi = N$ exactly. Given two factorizations
$[f_1]\cdots[f_k]$, $[f_1']\cdots[f_\ell']$, let $F : I\times I \to X$ be a homotopy between the
composed paths. Partition $I\times I$ into rectangles $R_{ij}$ each mapped by $F$ into a
single $A_{ij}$, refining so each point lies in at most three rectangles, and relabel them
$R_1,\ldots,R_{mn}$ in order. For a path $\gamma$ crossing $I\times I$ left to right, $F|\gamma$ is a
loop at $x_0$; let $\gamma_r$ separate the first $r$ rectangles from the rest, so $\gamma_0,\gamma_{mn}$
are the bottom/top edges. At each vertex $v$ of the grid with $F(v)\ne x_0$, choose a
path $g_v$ from $x_0$ to $F(v)$ inside the (path-connected, by hypothesis) intersection
of the $A_{ij}$'s meeting at $v$; inserting $\bar g_vg_v$ at vertices gives a factorization of
$[F|\gamma_r]$ for each $r$, well-defined up to equivalence, and consecutive $\gamma_r,\gamma_{r+1}$
give equivalent factorizations since $F|\gamma_r \simeq F|\gamma_{r+1}$ within a single $A_{ij}$. Choosing
the $g_v$'s along the bottom and top edges compatibly with the given $f_i$'s and
$f_i'$'s shows the factorization of $\gamma_0$ is equivalent to $[f_1]\cdots[f_k]$ and that of
$\gamma_{mn}$ to $[f_1']\cdots[f_\ell']$, so these two are equivalent to each other.
\end{proof}

\begin{example}[Linking of circles, via van Kampen]
\label{ex:linking-of-circles-fundamental-groups}
\dcref{ch1:1.23}
\group{Fundamental Group}
\uses{thm:van-kampen-theorem,prop:pi1-of-product-space,prop:pi1-sphere-trivial-n-geq-2,ex:borromean-rings-commutator}
We can now compute the fundamental groups of the three complements from
\cref{ex:borromean-rings-commutator}. The complement $\R^3-A$ of a single circle
deformation retracts onto $S^1\vee S^2$, so $\pi_1(\R^3-A) \cong \pi_1(S^1\vee S^2) \cong \Z$
since $\pi_1(S^2)=0$ (\cref{prop:pi1-sphere-trivial-n-geq-2}). The complement
$\R^3-(A\cup B)$ of two \emph{unlinked} circles deformation retracts onto
$S^1\vee S^1\vee S^2\vee S^2$, giving $\pi_1 \cong \Z*\Z$. The complement of two
\emph{linked} circles deformation retracts onto $S^2 \vee (S^1\times S^1)$ (a sphere
wedged with a separating torus), giving
$\pi_1 \cong \pi_1(S^1\times S^1) \cong \Z\times\Z$ (\cref{prop:pi1-of-product-space}) --
matching the informal claims made at the start of the chapter.


\end{example}

\begin{example}[Torus knot complements]
\label{ex:torus-knot-complement-fundamental-group}
\dcref{ch1:1.24}
\group{Fundamental Group}
\uses{thm:van-kampen-theorem,def:mapping-cylinder}
For relatively prime positive integers $m,n$, the \textbf{torus knot} $K=K_{m,n}\subset\R^3$
is the image of the embedding $f(z)=(z^m,z^n) : S^1 \to S^1\times S^1 \subset \R^3$, winding
$m$ times longitudinally and $n$ times meridionally.



Passing to the one-point compactification $S^3$ does not change $\pi_1$ (by van
Kampen, writing $S^3-K$ as $\R^3-K$ union a simply-connected neighborhood of
infinity meeting $\R^3-K$ in a simply-connected set $\approx S^2\times\R$). One shows
$S^3-K$ deformation retracts onto a complex $X=X_{m,n}$, the quotient of
$S^1\times I$ identifying $(z,0)\sim(e^{2\pi i/m}z,0)$ and $(z,1)\sim(e^{2\pi i/n}z,1)$; its two
halves $X_m,X_n$ (mapping cylinders of $z\mapsto z^m$, $z\mapsto z^n$,
\cref{def:mapping-cylinder}) meet in a circle $X_m\cap X_n$. All three pieces have
fundamental group $\Z$, and van Kampen's theorem gives
$\pi_1(X) = G_{m,n} := \langle a,b \mid a^m=b^n\rangle$, generated by the images of
generators of $\pi_1(X_m),\pi_1(X_n)$, with the relation coming from the common
generator of $\pi_1(X_m\cap X_n)$ mapping to $m$ times a generator in $X_m$ and $n$
times a generator in $X_n$.

If $m=1$ or $n=1$, $G_{m,n} \cong \Z$. For $m,n>1$: the element $a^m=b^n$
generates a central cyclic subgroup $C$ (it commutes with both generators), and
$G_{m,n}/C \cong \Z_m * \Z_n$; since a free product of nontrivial groups has trivial
center (Exercise 1 below) and all its torsion is conjugate into $\Z_m$ or $\Z_n$, the
integers $m,n$ are recovered (up to order) from $\Z_m*\Z_n$ alone -- via the order
$mn$ of its abelianization $\Z_m\times\Z_n$ together with the maximum order of a
torsion element, $\max(m,n)$ -- hence also from $G_{m,n}$, and $C$ is exactly the
center of $G_{m,n}$. (As shown in \cref{ex:gmn-action-on-tree}, $a$ and $b$ in fact
have infinite order in $G_{m,n}$, so $C$ is infinite cyclic, though this is not needed here.)

The construction of $X_{m,n}$ did not need $\gcd(m,n)=1$; that hypothesis is only
needed to realize $X_{m,n}$ as an actual knot complement, and in fact $X_{m,n}$
embeds in $\R^3$ only when $m,n$ are coprime (shown in the remarks after
Corollary 3.46, not yet transcribed) -- generalizing the nonembeddability of the
Klein bottle $X_{2,2}$ in $\R^3$.
\end{example}

\begin{example}[The shrinking wedge of circles]
\label{ex:shrinking-wedge-circles}
\dcref{ch1:1.25}
\group{Fundamental Group}
\uses{thm:fundamental-group-of-circle}
Let $X\subset\R^2$ be the union of circles $C_n$ of radius $1/n$ centered at $(1/n,0)$,
$n=1,2,\ldots$. Although superficially like a wedge of infinitely many circles, $X$ has
a much larger fundamental group: the retractions $r_n:X\to C_n$ (collapsing all
other $C_i$ to the origin) induce a surjective homomorphism
$\rho:\pi_1(X)\to\prod_\infty\Z$ (the full direct product, not the direct sum), since any
sequence of integers $(k_n)$ is realized by a loop winding $k_n$ times around $C_n$
during $[1-1/n,1-1/(n+1)]$ -- continuous at time $1$ because every neighborhood
of the basepoint contains all but finitely many $C_n$. Hence $\pi_1(X)$ is uncountable,
unlike the countable fundamental group of a genuine wedge of countably many
circles. It is also nonabelian: collapsing all circles smaller than $C_n$ surjects
$\pi_1(X)$ onto a free group of rank $n$, for every $n$. (For a complete description see
\cite{CannonConner2000}; it is a theorem of \cite{Shelah1988}
that for path-connected, locally path-connected compact metric $X$, $\pi_1(X)$ is
either finitely generated or uncountable.)


\end{example}

\subsection*{Applications to Cell Complexes}

For the rest of this section we consider how $\pi_1$ is affected by attaching 2-cells.
Suppose 2-cells $e_\alpha^2$ are attached to a path-connected space $X$ via maps
$\varphi_\alpha : S^1 \to X$, producing $Y$. Choosing a basepoint $x_0\in X$ and a path $y_\alpha$
from $x_0$ to $\varphi_\alpha(s_0)$, the loop $y_\alpha\varphi_\alpha\bar y_\alpha$ at $x_0$ becomes nullhomotopic
once $e_\alpha^2$ is attached, so the normal subgroup $N \subset \pi_1(X,x_0)$ generated by all
these loops lies in $\ker(\pi_1(X,x_0)\to\pi_1(Y,x_0))$.

\begin{proposition}[Attaching cells and $\pi_1$]
\label{prop:attaching-cells-pi1}
\dcref{ch1:1.26}
\group{Fundamental Group}
\uses{thm:van-kampen-theorem,prop:pi1-sphere-trivial-n-geq-2}
\begin{enumerate}
\item[(a)] If $Y$ is obtained from $X$ by attaching 2-cells as above, the inclusion
$X\hookrightarrow Y$ induces a surjection $\pi_1(X,x_0)\to\pi_1(Y,x_0)$ with kernel exactly $N$,
so $\pi_1(Y) \approx \pi_1(X)/N$.
\item[(b)] If $Y$ is obtained from $X$ by attaching $n$-cells for fixed $n>2$, the
inclusion induces an isomorphism $\pi_1(X,x_0)\approx\pi_1(Y,x_0)$.
\item[(c)] For a path-connected cell complex $X$, the inclusion of the 2-skeleton
$X^2\hookrightarrow X$ induces an isomorphism $\pi_1(X^2,x_0)\approx\pi_1(X,x_0)$.
\end{enumerate}
\end{proposition}

(Part (a) shows in particular that $N$ does not depend on the choice of paths
$y_\alpha$: replacing $y_\alpha$ by $\eta_\alpha$ with the same endpoints conjugates
$y_\alpha\varphi_\alpha\bar y_\alpha$ by $\eta_\alpha\bar y_\alpha$.)

\begin{proof}
(a) Expand $Y$ to $Z$, obtained by attaching rectangular strips $S_\alpha=I\times I$ with
lower edge along $y_\alpha$, right edge along a radial arc into $e_\alpha^2$ from $\varphi_\alpha(s_0)$, and
all left edges identified together; $Z$ deformation retracts onto $Y$ (the strips' top
edges are free). Choosing a point $y_\alpha$ in each $e_\alpha^2$ off the attaching arc, let
$A = Z-\bigcup_\alpha\{y_\alpha\}$ (deformation retracts onto $X$) and $B = Z-X$
(contractible). Van Kampen's theorem applied to $\{A,B\}$ gives $\pi_1(Z)$ as the
quotient of $\pi_1(A)$ by the normal subgroup generated by the image of
$\pi_1(A\cap B)\to\pi_1(A)$, and a second application of van Kampen's theorem (to the
cover of $A\cap B$ by $A_\alpha = A\cap B - \bigcup_{\beta\ne\alpha}e_\beta^2$, each $\pi_1(A_\alpha)\cong\Z$) shows
$\pi_1(A\cap B)$ is generated by loops $\delta_\alpha$ corresponding to the
$[y_\alpha\varphi_\alpha\bar y_\alpha]$'s, giving (a).

(b) Same argument with $e_\alpha^n$: now $A_\alpha$ deformation retracts onto $S^{n-1}$, so
$\pi_1(A_\alpha)=0$ for $n>2$ by \cref{prop:pi1-sphere-trivial-n-geq-2}, hence $\pi_1(A\cap B)=0$.

(c) Follows from (b) by induction when $X$ is finite-dimensional. In general, a loop
$f$ at $x_0\in X^2$ has compact image in some $X^n$ (Proposition A.1 in the
Appendix, not yet transcribed), so (b) gives $f$ homotopic to a loop in $X^2$,
proving surjectivity of $\pi_1(X^2,x_0)\to\pi_1(X,x_0)$; injectivity is similar, using
that a nullhomotopy also has compact image in some $X^n$.
\end{proof}



As a first application, the closed orientable surface $M_g$ of genus $g$ (one 0-cell,
$2g$ 1-cells, one 2-cell, \cref{sec:cell-complexes}) has 1-skeleton a wedge of $2g$ circles
(free fundamental group on $2g$ generators), with the 2-cell attached along the
product of commutators $[a_1,b_1]\cdots[a_g,b_g]$, so
\[
\pi_1(M_g) \cong \langle a_1,b_1,\ldots,a_g,b_g \mid [a_1,b_1]\cdots[a_g,b_g]\rangle,
\]
using $\langle g_\alpha \mid r_\beta\rangle$ for the free group on the $g_\alpha$'s modulo the normal
subgroup generated by the words $r_\beta$.

\begin{corollary}[Genus distinguishes closed orientable surfaces]
\label{cor:genus-surfaces-distinct}
\dcref{ch1:1.27}
\group{Fundamental Group}
\uses{prop:attaching-cells-pi1}
$M_g$ is not homeomorphic, or even homotopy equivalent, to $M_h$ if $g\ne h$.
\end{corollary}

\begin{proof}
The abelianization of $\pi_1(M_g)$ is $\Z^{2g}$. If $M_g\simeq M_h$ then
$\pi_1(M_g)\cong\pi_1(M_h)$, so their abelianizations agree, forcing $g=h$.
\end{proof}

Nonorientable surfaces $N_g$ are obtained by attaching a 2-cell to a wedge of $g$
circles via the word $a_1^2\cdots a_g^2$; $N_1 = \RP^2$ and $N_2$ is the Klein bottle
(usually presented instead as a square with sides identified via $aba^{-1}b$, or, after
cutting along a diagonal, via $a^2c^2$). By \cref{prop:attaching-cells-pi1},
$\pi_1(N_g) \cong \langle a_1,\ldots,a_g \mid a_1^2\cdots a_g^2\rangle$, which abelianizes to
$\Z_2 \oplus \Z^{g-1}$ (rechoosing generators as $a_1,\ldots,a_{g-1}$ and
$a_1+\cdots+a_g$). Hence $N_g \not\simeq N_h$ for $g\ne h$, and no $N_g$ is homotopy
equivalent to any $M_h$.



\begin{corollary}[Every group is $\pi_1$ of some 2-complex]
\label{cor:every-group-realized-as-pi1}
\dcref{ch1:1.28}
\group{Fundamental Group}
\uses{prop:attaching-cells-pi1,ex:wedge-sum-fundamental-group}
For every group $G$ there is a 2-dimensional cell complex $X_G$ with
$\pi_1(X_G) \cong G$.
\end{corollary}

\begin{proof}
Choose a presentation $G = \langle g_\alpha \mid r_\beta\rangle$ (every group is a quotient of a
free group). Build $X_G$ from $\bigvee_\alpha S^1_\alpha$ (\cref{ex:wedge-sum-fundamental-group}) by
attaching 2-cells $e_\beta^2$ along the loops specified by the words $r_\beta$.
\end{proof}

\begin{example}[$X_G$ for $G=\Z_n$]
\label{ex:cyclic-group-complex-xg}
\dcref{ch1:1.29}
\group{Fundamental Group}
\uses{cor:every-group-realized-as-pi1}
For $G = \langle a \mid a^n\rangle = \Z_n$, $X_G$ is $S^1$ with a 2-cell attached via
$z\mapsto z^n$ ($S^1\subset\C$). For $n=2$, $X_G = \RP^2$; for $n>2$, $X_G$ is not a
surface, since $n$ `sheets' of the 2-cell meet along $S^1$. For instance at $n=3$ a
neighborhood of $S^1$ in $X_G$ is built from the graph $\Upsilon$ (three edges from a
point) crossed with $I$, glued end to end with a one-third twist; its boundary is a
single circle, to which the rest of the 2-cell is attached -- a construction that needs
$\R^4$, not $\R^3$.


\end{example}

\subsection*{Exercises}
% Exercises are transcribed for completeness but are not given blueprint labels or
% \uses edges.

\begin{enumerate}
\item Show that the free product $G*H$ of nontrivial groups $G,H$ has trivial
center, and that its only finite-order elements are conjugates of finite-order
elements of $G$ or $H$.

\item Let $X\subset\R^m$ be the union of convex open sets $X_1,\ldots,X_n$ with
$X_i\cap X_j\cap X_k=\emptyset$ for all $i,j,k$. Show $X$ is simply-connected.

\item Show that the complement of a finite set of points in $\R^n$ is
simply-connected if $n\ge 3$.

\item Let $X\subset\R^3$ be the union of $n$ lines through the origin. Compute
$\pi_1(\R^3-X)$.

\item Let $X\subset\R^2$ be a connected graph that is a finite union of straight line
segments. Show $\pi_1(X)$ is free with basis the boundaries of the bounded
complementary regions, joined to a basepoint by suitable paths in $X$. [Assume
the Jordan curve theorem for polygonal simple closed curves.]

\item Use \cref{prop:attaching-cells-pi1} to show that the complement of a closed
discrete subspace of $\R^n$ is simply-connected if $n\ge 3$.

\item Let $X$ be $S^2$ with north and south poles identified. Put a cell structure
on $X$ and compute $\pi_1(X)$.

\item Compute the fundamental group of the space obtained from two tori
$S^1\times S^1$ by identifying a circle $S^1\times\{x_0\}$ in one with the
corresponding circle in the other.

\item In the surface $M_g$, let $C$ be a circle separating $M_g$ into subsurfaces
$M_h'$ and $M_k'$ (from $M_h$, $M_k$ minus a disk each). Show $M_h'$ does not
retract onto $C$, hence $M_g$ does not retract onto $C$ [abelianize $\pi_1$], but
$M_g$ does retract onto a nonseparating circle $C'$.

\item Two arcs $\alpha,\beta$ are embedded in $D^2\times I$ with a loop $\gamma$
nullhomotopic in $D^2\times I$; show $\gamma$ is not nullhomotopic in the
complement of $\alpha\cup\beta$.

\item The \textbf{mapping torus} $T_f$ of $f:X\to X$ is $X\times I$ with $(x,0)$
identified to $(f(x),1)$. For $X=S^1\vee S^1$ (or $X=S^1\times S^1$) and $f$
basepoint-preserving, compute a presentation of $\pi_1(T_f)$ in terms of
$f_*:\pi_1(X)\to\pi_1(X)$. [Regard $T_f$ as built from $X\vee S^1$ by attaching
cells.]

\item Let $X\subset\R^3$ be the usual Klein-bottle-with-self-intersection picture,
and $Y\subset X$ the result of deleting the open disk bounded by the
self-intersection circle. Show $\pi_1(X)\cong\Z\times\Z$ and $\pi_1(Y)$ has
presentation $\langle a,b,c \mid aba^{-1}b^{-1}cbc^{-1}\rangle$ up to a sign, and that
$\pi_1(Y) \cong \pi_1(\R^3-Z)$ for a certain graph $Z$. ($\pi_1(X)$ and $\pi_1(Y)$ are
in fact non-isomorphic, but this is not easy to show; see Example 1B.13, not yet
transcribed.)

\item The space $Y$ above can also be obtained from a disk with two holes by
identifying its three boundary circles in one of two essentially different ways;
show the other way gives $Z$ with $\pi_1(Z) \not\cong \pi_1(Y)$ [abelianize].

\item Consider the cube $I^3$ with each square face identified to the opposite
face by a translation plus a one-quarter twist. Show the quotient $X$ is a cell
complex with two 0-cells, four 1-cells, three 2-cells, one 3-cell, and that
$\pi_1(X)$ is the quaternion group $\{\pm1,\pm i,\pm j,\pm k\}$ of order eight.

\item Given $(X,x_0)$, build a CW complex $L(X)$ with one 0-cell, a 1-cell
$e_\gamma^1$ for each loop $\gamma$ at $x_0$, and a 2-cell $e_\tau^2$ for each map $\tau$ of a
triangle $PQR$ to $X$ sending vertices to $x_0$, attached along the three loops
$\tau|_{PQ},\tau|_{PR},\tau|_{QR}$. Show the natural map $L(X)\to X$ induces an
isomorphism on $\pi_1$.

\item Show the fundamental group of the infinite-genus surface (infinitely many
torus handles attached in sequence) is free on infinitely many generators.

\item Show $\pi_1(\R^2-\Q^2)$ is uncountable.

\item Let $X\subset\R$ be $\{1,1/2,1/3,\ldots\}\cup\{0\}$. (a) For the suspension $SX$,
show $\pi_1(SX)$ is free on countably many generators, hence countable, unlike the
reduced suspension $\Sigma X$ (the shrinking wedge of \cref{ex:shrinking-wedge-circles}),
which has uncountable $\pi_1$. (b) Let $C$ be the mapping cone of $SX\to\Sigma X$
(so $C$ is the reduced suspension of the cone $CX$). Show $\pi_1(C)$ is uncountable
by constructing a surjection $\pi_1(C) \to \prod_\omega\Z/\bigoplus_\omega\Z$. Thus the
reduced suspension of a contractible space need not be contractible, unlike the
unreduced suspension.

\item Show the union of spheres of radius $1/n$ centered at $(1/n,0,0)$ in $\R^3$,
$n=1,2,\ldots$, is simply-connected.

\item Let $X\subset\R^2$ be the union of circles $C_n$ of radius $n$ centered at
$(n,0)$. Show $\pi_1(X) \cong *_n\pi_1(C_n)$, the same as for an infinite wedge
$\bigvee_\infty S^1$, and that $X \simeq \bigvee_\infty S^1$ though they are not
homeomorphic.

\item Show the join $X*Y$ of nonempty spaces is simply-connected if $X$ is
path-connected.

\item [\textbf{Wirtinger presentation.}] With a knot $K\subset\R^3$ positioned to lie
almost flat with finitely many crossing arcs $\beta_\ell$ over table-top arcs $\alpha_i$, build
a 2-complex $X$ deformation-retracting onto $\R^3-K$ from a rectangle $T$ (the
table), a strip $R_i$ over each $\alpha_i$, and a square $S_\ell$ over each $\beta_\ell$ glued to three
strips. (a) Show $\pi_1(\R^3-K)$ has a presentation with one generator $x_i$ per strip
and one relation $x_ix_jx_i^{-1}=x_k$ per square $S_\ell$. (b) Use this to show the
abelianization of $\pi_1(\R^3-K)$ is $\Z$.
\end{enumerate}

\section{1.3 Covering Spaces}
\label{sec:covering-spaces}

We come now to the second main topic of this chapter, covering spaces, already
encountered briefly via the helix projection $\R \to S^1$
(\cref{thm:fundamental-group-of-circle}). One of the main results of this section is an exact
correspondence between connected covering spaces of $X$ and subgroups of
$\pi_1(X)$ -- strikingly reminiscent of the Galois correspondence between field
extensions and subgroups of the Galois group.

\begin{definition}[Sheets of a covering space]
\label{def:covering-space-sheets}
\group{Covering Spaces}
\uses{def:covering-space}
For an evenly covered open $U \subset X$ (\cref{def:covering-space}), the disjoint open
sets of $\tilde X$ projecting homeomorphically to $U$ are called the \textbf{sheets} of
$\tilde X$ over $U$. If $U$ is connected the sheets are exactly the components of
$p^{-1}(U)$. The \textbf{number of sheets} over $U$ is $|p^{-1}(x)|$ for $x \in U$; this is
locally constant in $x$, hence constant if $X$ is connected.
\end{definition}

An example is the helicoid $S = \{(s\cos 2\pi t, s\sin 2\pi t, t) : (s,t) \in (0,\infty)\times\R\}$,
projecting onto $\R^2-\{0\}$ by $(x,y,z)\mapsto(x,y)$, a covering space with
countably infinitely many sheets over each small disk. Another is
$p : S^1 \to S^1$, $p(z) = z^n$, an $n$-sheeted covering. Together with the helix
(infinite-sheeted), these exhaust the connected covering spaces of $S^1$.

The covering spaces of $X = S^1\vee S^1$ form a rich illustrative family: viewing $X$
as a graph with one vertex and two labeled, oriented edges $a,b$, a covering space
corresponds to a \textbf{2-oriented graph} $\tilde X$ (four edge-ends at each vertex,
labeled/oriented to match $X$ locally), with $p$ sending vertices to the vertex and
edges to same-labeled edges preserving orientation. Every graph with four edge-ends
per vertex can be 2-oriented (via an Eulerian circuit argument for finite graphs; see
\cite{Konig1990} for the infinite case).



A simply-connected covering space of $S^1\vee S^1$ is built by iteratively adjoining
perpendicular open segments of geometrically shrinking length $2\lambda^{n-1}$ at
distance $\lambda^{n-1}$ from all previous endpoints (for fixed $0<\lambda<1/2$),
labeling horizontal edges $a$ and vertical edges $b$; this is called the
\textbf{universal cover} of $X$ because (as the general theory below shows) it covers
every other connected covering space of $X$.



The covering spaces (1)--(14) above are all non-simply-connected, with free
fundamental group on the listed words in $a,b$ (provable by van Kampen's theorem),
which also generate the image subgroup $p_*(\pi_1(\tilde X,\tilde x_0))$ in
$\pi_1(X,x_0)=\langle a,b\rangle$; since $p_*$ is always injective
(\cref{prop:covering-map-pi1-injective}), this shows $\Z*\Z$ contains free subgroups of
every finite rank, and even countably infinite rank (examples (10), (11)). Changing
the basepoint vertex conjugates the subgroup by the projection of a connecting path
in $\tilde X$; e.g.\ (3) and (4) differ only by basepoint, with subgroups conjugate by
$b$. An \textbf{isomorphism} of covering spaces of $X$ is a graph isomorphism
preserving edge labels and orientations; (3),(4) are isomorphic (not basepoint-
preservingly), while (5),(6) have homeomorphic but non-isomorphic underlying
graphs, so their subgroups are isomorphic but not conjugate. The most symmetric
covering spaces -- (1),(2),(5)--(8),(11) -- are exactly the \textbf{normal} ones
(\cref{def:normal-covering-space}), where the corresponding subgroup is normal in
$\pi_1(X,x_0)$, with symmetry group $\pi_1(X,x_0)/(\text{subgroup})$; since every
two-generator group is a quotient of $\Z*\Z$, every two-generator group arises this
way as a symmetry group of some covering space of $X$.

\subsection*{Lifting Properties}

\begin{definition}[Lift of a map]
\label{def:lift-of-a-map}
\group{Covering Spaces}
\uses{def:covering-space}
Given $p : \tilde X \to X$, a \textbf{lift} of $f : Y \to X$ is a map $\tilde f : Y \to \tilde X$ with
$p\tilde f = f$.
\end{definition}

\begin{proposition}[Homotopy lifting property]
\label{prop:covering-homotopy-property}
\dcref{ch1:1.30}
\group{Covering Spaces}
\uses{def:covering-space,def:lift-of-a-map,thm:fundamental-group-of-circle}
Given a covering space $p : \tilde X \to X$, a homotopy $f_t : Y \to X$, and a lift
$\tilde f_0 : Y \to \tilde X$ of $f_0$, there is a unique homotopy $\tilde f_t : Y \to \tilde X$ of
$\tilde f_0$ lifting $f_t$.
\end{proposition}

This was already proved as step (c) in the proof of
\cref{thm:fundamental-group-of-circle}. Taking $Y$ a point gives \textbf{path lifting}: every
path lifts uniquely from any chosen starting lift of its initial point. Taking $Y=I$,
every homotopy of a path lifts to a homotopy of paths (endpoints fixed, since each
endpoint traces a lift of a constant path, hence is constant).

\begin{proposition}[Covering maps are injective on $\pi_1$]
\label{prop:covering-map-pi1-injective}
\dcref{ch1:1.31}
\group{Covering Spaces}
\uses{prop:covering-homotopy-property,def:loop-and-pi1-set}
The map $p_* : \pi_1(\tilde X,\tilde x_0) \to \pi_1(X,x_0)$ induced by a covering space
$p : (\tilde X,\tilde x_0) \to (X,x_0)$ is injective, with image the classes of loops at $x_0$
whose lift to $\tilde X$ starting at $\tilde x_0$ is a loop.
\end{proposition}

\begin{proof}
If $\tilde f_0$ is a loop with $[p\tilde f_0]=0$ via a homotopy $f_t$ to the trivial loop, the
lifted homotopy (by \cref{prop:covering-homotopy-property}) goes from $\tilde f_0$ to a
constant loop, so $[\tilde f_0]=0$. For the image description: loops lifting to loops
clearly lie in the image; conversely a loop representing an image class is homotopic
to one with such a lift, so by homotopy lifting it has one too.
\end{proof}

\begin{proposition}[Sheets equal the index of the image subgroup]
\label{prop:covering-sheets-equals-index}
\dcref{ch1:1.32}
\group{Covering Spaces}
\uses{def:covering-space-sheets,prop:covering-map-pi1-injective}
For $p:(\tilde X,\tilde x_0)\to(X,x_0)$ with $X,\tilde X$ path-connected, the number of
sheets equals the index of $H = p_*(\pi_1(\tilde X,\tilde x_0))$ in $\pi_1(X,x_0)$.
\end{proposition}

\begin{proof}
Send the coset $H[g]$ to $\tilde g(1)$, the endpoint of the lift of $g$ starting at
$\tilde x_0$. This is well-defined and surjective by path-connectedness of $\tilde X$; it is
injective since $\Phi(H[g_1])=\Phi(H[g_2])$ means $g_1\bar g_2$ lifts to a loop, so
$[g_1][g_2]^{-1}\in H$.
\end{proof}

\begin{proposition}[Lifting criterion]
\label{prop:covering-lifting-criterion}
\dcref{ch1:1.33}
\group{Covering Spaces}
\uses{prop:covering-homotopy-property,def:lift-of-a-map}
For $p:(\tilde X,\tilde x_0)\to(X,x_0)$ and $f:(Y,y_0)\to(X,x_0)$ with $Y$
path-connected and locally path-connected, a lift $\tilde f:(Y,y_0)\to(\tilde X,\tilde x_0)$
exists iff $f_*(\pi_1(Y,y_0)) \subset p_*(\pi_1(\tilde X,\tilde x_0))$.
\end{proposition}

\begin{proof}
Necessity is clear from $f_*=p_*\tilde f_*$. For sufficiency, define $\tilde f(y)=\widetilde{f\gamma}(1)$
for any path $\gamma$ from $y_0$ to $y$, lifted starting at $\tilde x_0$; this is independent
of $\gamma$ by the subgroup hypothesis (via homotopy lifting applied to the loop
$(f\gamma')\cdot\overline{f\gamma}$), and continuous by locally choosing $\gamma$ through a
path-connected neighborhood $V$ with $f(V)$ inside an evenly covered set.
\end{proof}

Local path-connectedness of $Y$ is genuinely needed (Exercise 7 below gives a
counterexample).

\begin{proposition}[Unique lifting property]
\label{prop:covering-unique-lifting}
\dcref{ch1:1.34}
\group{Covering Spaces}
\uses{def:lift-of-a-map}
Given $p:\tilde X\to X$ and $f:Y\to X$ with $Y$ connected, if two lifts
$\tilde f_1,\tilde f_2$ agree at one point, they agree everywhere.
\end{proposition}

\begin{proof}
The set where $\tilde f_1=\tilde f_2$ is open (they agree on a whole sheet if they agree
at a point, by injectivity of $p$ on that sheet) and closed (if they disagree at $y$,
their distinct containing sheets are disjoint, so they disagree on a neighborhood),
hence all of connected $Y$ or empty at that point.
\end{proof}

\subsection*{The Classification of Covering Spaces}

We now restrict attention to $X$ locally path-connected (so path-components =
components) and, for classification purposes, connected; connected covering spaces
of such $X$ are then the same as path-connected ones.

\begin{definition}[Semilocally simply-connected]
\label{def:semilocally-simply-connected}
\group{Covering Spaces}
$X$ is \textbf{semilocally simply-connected} if each $x\in X$ has a neighborhood $U$
with $\pi_1(U,x)\to\pi_1(X,x)$ trivial.
\end{definition}

This is necessary for $X$ to have a simply-connected covering space (a nullhomotopy
of a lifted loop in $\tilde U$ projects to one in $X$), and holds automatically for
locally simply-connected spaces such as CW complexes (locally contractible, shown
in the Appendix). The shrinking wedge of circles (\cref{ex:shrinking-wedge-circles}) is
not semilocally simply-connected, though its cone is (being contractible, though not
locally simply-connected).

\begin{proposition}[Construction of the universal cover]
\label{prop:universal-cover-construction}
\group{Covering Spaces}
\uses{def:semilocally-simply-connected,prop:simply-connected-iff-unique-path-class}
If $X$ is path-connected, locally path-connected, and semilocally simply-connected,
then $X$ has a simply-connected covering space.
\end{proposition}

\begin{proof}
Set $\tilde X = \{[\gamma] : \gamma \text{ a path in } X \text{ starting at } x_0\}$ and
$p[\gamma]=\gamma(1)$, surjective since $X$ is path-connected. Let $\mathcal U$ be the
path-connected open $U\subset X$ with $\pi_1(U)\to\pi_1(X)$ trivial (a basis, given the
hypotheses); for $U\in\mathcal U$ and a path $\gamma$ to a point of $U$, set
$U_{[\gamma]} = \{[\gamma\cdot\eta] : \eta \text{ a path in } U \text{ from } \gamma(1)\}$, which depends
only on $[\gamma]$ and satisfies $U_{[\gamma]}=U_{[\gamma']}$ whenever $[\gamma']\in U_{[\gamma]}$
(since $\pi_1(U)\to\pi_1(X)$ is trivial). These properties make the $U_{[\gamma]}$ a basis
for a topology on $\tilde X$, and $p:U_{[\gamma]}\to U$ a homeomorphism for each,
matching sub-basic opens $V_{[\gamma']}\subset U_{[\gamma]}$ with $V\in\mathcal U$, $V\subset U$; hence
$p$ is a covering map, the sets $U_{[\gamma]}$ partitioning $p^{-1}(U)$.

For $[\gamma]\in\tilde X$, the path $t\mapsto[\gamma_t]$ (where $\gamma_t$ is $\gamma$ stopped at time
$t$) lifts $\gamma$ from $[x_0]$ (the constant path) to $[\gamma]$, so $\tilde X$ is
path-connected. If a loop $\gamma$ at $x_0$ lifts to a loop at $[x_0]$ then
$[\gamma_1]=[\gamma]=[x_0]$, i.e.\ $\gamma$ is nullhomotopic; since $p_*$ is injective
(\cref{prop:covering-map-pi1-injective}), this shows $\pi_1(\tilde X,[x_0])=0$.
\end{proof}

\begin{example}[Universal cover of $X_{m,n}$ by gluing]
\label{ex:xmn-universal-cover-construction}
\dcref{ch1:1.35}
\group{Covering Spaces}
\uses{ex:torus-knot-complement-fundamental-group,def:covering-space}
For $m,n\ge 2$ let $X_{m,n}$ be as in \cref{ex:torus-knot-complement-fundamental-group},
the union of mapping cylinders $A,B$ of $z\mapsto z^m$, $z\mapsto z^n$ glued along a
common circle $A\cap B=S^1$ (so $X_{2,2}$ is the Klein bottle). The universal cover
$\tilde A \cong C_m\times\R$ (with $C_m$ the cone on $m$ points, a $m$-pronged star
graph), similarly $\tilde B\cong C_n\times\R$. Gluing infinitely many copies of $\tilde A,\tilde B$
along their free boundary lines in an alternating tree pattern builds
$\tilde X_{m,n} = T_{m,n}\times\R$, where $T_{m,n}$ is the analogous infinite tree of
copies of $C_m,C_n$. Each finite stage of $T_{m,n}$ deformation retracts onto the
previous one, so (concatenating these over shrinking time intervals) $T_{m,n}$
deformation retracts onto a single $C_m$, hence is contractible; so is
$\tilde X_{m,n} = T_{m,n}\times\R$, in particular simply-connected. The evident map
sending each copy of $\tilde A,\tilde B$ to $A,B$ is then the universal covering
$\tilde X_{m,n}\to X_{m,n}$.


\end{example}

\begin{proposition}[Covering spaces realizing arbitrary subgroups]
\label{prop:covering-space-existence-for-subgroup}
\dcref{ch1:1.36}
\group{Covering Spaces}
\uses{prop:universal-cover-construction}
If $X$ is path-connected, locally path-connected, and semilocally simply-connected,
then for every subgroup $H\subset\pi_1(X,x_0)$ there is a covering space
$p:X_H\to X$ with $p_*(\pi_1(X_H,\tilde x_0))=H$ for a suitable basepoint $\tilde x_0$.
\end{proposition}

\begin{proof}
On the universal cover $\tilde X$, define $[\gamma]\sim[\gamma']$ iff $\gamma(1)=\gamma'(1)$ and
$[\gamma\cdot\bar\gamma']\in H$; this is an equivalence relation since $H$ is a subgroup.
Let $X_H = \tilde X/{\sim}$. Since $[\gamma]\sim[\gamma']$ (with matching endpoints) implies
$[\gamma\eta]\sim[\gamma'\eta]$ for any further path $\eta$, whole basic neighborhoods
$U_{[\gamma]}$ get identified together whenever a point of each does, so the induced map
$X_H\to X$ is a covering space. Taking $\tilde x_0$ to be the class of the constant path,
a loop $\gamma$ at $x_0$ lifts to a loop in $X_H$ iff $[\gamma]\sim[x_0]$, i.e.\ $[\gamma]\in H$, so
$p_*(\pi_1(X_H,\tilde x_0)) = H$.
\end{proof}

\begin{proposition}[Basepointed covering spaces are isomorphic iff same image subgroup]
\label{prop:covering-space-isomorphism-criterion}
\dcref{ch1:1.37}
\group{Covering Spaces}
\uses{prop:covering-lifting-criterion,prop:covering-unique-lifting}
If $X$ is path-connected and locally path-connected, path-connected covering
spaces $p_1:\tilde X_1\to X$, $p_2:\tilde X_2\to X$ admit a basepoint-preserving
isomorphism $\tilde X_1\to\tilde X_2$ iff
$p_{1*}(\pi_1(\tilde X_1,\tilde x_1)) = p_{2*}(\pi_1(\tilde X_2,\tilde x_2))$.
\end{proposition}

\begin{proof}
Necessity is immediate from $p_1=p_2f$. For sufficiency, the lifting criterion
(\cref{prop:covering-lifting-criterion}) gives basepoint-preserving lifts
$\tilde p_1:\tilde X_1\to\tilde X_2$ of $p_1$ and $\tilde p_2:\tilde X_2\to\tilde X_1$ of $p_2$; by
uniqueness of lifts (\cref{prop:covering-unique-lifting}), their compositions (which fix
basepoints and lift the identity) are the identity, so they are inverse isomorphisms.
\end{proof}

\begin{theorem}[Classification of covering spaces]
\label{thm:covering-space-classification}
\dcref{ch1:1.38}
\group{Covering Spaces}
\uses{prop:covering-space-existence-for-subgroup,prop:covering-space-isomorphism-criterion}
Let $X$ be path-connected, locally path-connected, and semilocally
simply-connected. Associating $p_*(\pi_1(\tilde X,\tilde x_0))$ to $p:(\tilde X,\tilde x_0)\to(X,x_0)$
gives a bijection between basepoint-preserving isomorphism classes of
path-connected covering spaces of $X$ and subgroups of $\pi_1(X,x_0)$. Ignoring
basepoints, this gives a bijection between isomorphism classes of path-connected
covering spaces and conjugacy classes of subgroups of $\pi_1(X,x_0)$.
\end{theorem}

\begin{proof}
Existence and the basepointed bijection are
\cref{prop:covering-space-existence-for-subgroup,prop:covering-space-isomorphism-criterion}. For the
unbased statement: if $\tilde\gamma$ is a path from basepoint $\tilde x_0$ to another lift
$\tilde x_1$ of $x_0$, projecting to a loop $\gamma$ representing $g\in\pi_1(X,x_0)$, then
$H_1 := p_*(\pi_1(\tilde X,\tilde x_1)) = g^{-1}H_0g$ where $H_0 := p_*(\pi_1(\tilde X,\tilde x_0))$
(shown via $\bar\gamma\cdot f\cdot\gamma$ for loops $f$ at $\tilde x_0$, in both directions), and
conversely every conjugate of $H_0$ arises this way by choosing $\gamma$ representing
the appropriate $g$.
\end{proof}

\begin{definition}[Universal cover]
\label{def:universal-cover}
\group{Covering Spaces}
\uses{thm:covering-space-classification,prop:covering-lifting-criterion}
A simply-connected covering space of a path-connected, locally path-connected
$X$ is called \textbf{the universal cover} of $X$: it covers every other path-connected
covering space of $X$ (by the lifting criterion applied to the trivial subgroup), and is
unique up to isomorphism by \cref{thm:covering-space-classification}.
\end{definition}

More generally, path-connected covering spaces of $X$ are partially ordered by which
covers which, matching the partial ordering by inclusion of the corresponding
subgroups (or conjugacy classes) of $\pi_1(X)$.

\subsection*{Representing Covering Spaces by Permutations}

For a covering space $p:\tilde X\to X$, a path $\gamma$ from $x$ to $x'$ induces a bijection
$p^{-1}(x')\to p^{-1}(x)$ sending the endpoint of each lift to its start (this inverse
convention makes composition match: $L_{\gamma\cdot\eta}=L_\gamma L_\eta$). Restricting to loops
at $x_0$ gives a homomorphism $\pi_1(X,x_0)\to \mathrm{Sym}(p^{-1}(x_0))$, the
\textbf{action of $\pi_1(X,x_0)$ on the fiber}. Conversely, given any action $\rho$ of
$\pi_1(X,x_0)$ on a discrete set $F$, and a universal cover $\tilde X_0\to X$ (points =
homotopy classes of paths from $x_0$), the quotient $\tilde X_\rho$ of $\tilde X_0\times F$
by $([\gamma],\xi)\sim([\lambda\gamma],\rho(\lambda)\xi)$ for $[\lambda]\in\pi_1(X,x_0)$ carries a natural
covering map to $X$, and every covering space of $X$ arises this way (up to
isomorphism) from the action on its fiber over $x_0$; isomorphic covering spaces
correspond exactly to isomorphic $\pi_1(X,x_0)$-sets. Consequently, $n$-sheeted
covering spaces of $X$ are classified by conjugacy classes of homomorphisms
$\pi_1(X,x_0) \to \Sigma_n$.

\subsection*{Deck Transformations and Group Actions}

\begin{definition}[Deck transformation group]
\label{def:deck-transformation-group}
\group{Covering Spaces}
\uses{def:covering-space,prop:covering-unique-lifting}
For $p:\tilde X\to X$, the isomorphisms $\tilde X\to\tilde X$ (self-isomorphisms of the
covering space) are \textbf{deck transformations}, forming a group $G(\tilde X)$ under
composition. By unique lifting (\cref{prop:covering-unique-lifting}), a deck
transformation of path-connected $\tilde X$ is determined by its value at one point;
only the identity fixes a point.
\end{definition}

\begin{definition}[Normal covering space]
\label{def:normal-covering-space}
\group{Covering Spaces}
\uses{def:deck-transformation-group}
$p:\tilde X\to X$ is \textbf{normal} (or \textbf{regular}) if for every $x\in X$ and every
pair of lifts $\tilde x,\tilde x'$ there is a deck transformation taking $\tilde x$ to $\tilde x'$.
\end{definition}

\begin{proposition}[Deck group and normalizers]
\label{prop:normal-covering-deck-group}
\dcref{ch1:1.39}
\group{Covering Spaces}
\uses{def:normal-covering-space,thm:covering-space-classification,prop:covering-lifting-criterion}
Let $p:(\tilde X,\tilde x_0)\to(X,x_0)$ be a path-connected covering space of
path-connected, locally path-connected $X$, and $H = p_*(\pi_1(\tilde X,\tilde x_0))$. Then
(a) the covering is normal iff $H \trianglelefteq \pi_1(X,x_0)$; (b) $G(\tilde X) \cong N(H)/H$
for $N(H)$ the normalizer of $H$. In particular $G(\tilde X)\cong\pi_1(X,x_0)/H$ if
normal, and $G(\tilde X) \cong \pi_1(X)$ for the universal cover.
\end{proposition}

\begin{proof}
Changing basepoint $\tilde x_0\to\tilde x_1$ (both over $x_0$) conjugates $H$ by
$[\gamma]$ for $\gamma$ the projection of a connecting path (as in
\cref{thm:covering-space-classification}), so $[\gamma]\in N(H)$ iff
$p_*\pi_1(\tilde X,\tilde x_0) = p_*\pi_1(\tilde X,\tilde x_1)$, iff (by the lifting criterion) there is
a deck transformation $\tilde x_0\mapsto\tilde x_1$. Hence normality $\iff N(H)=\pi_1(X,x_0)
\iff H\trianglelefteq\pi_1(X,x_0)$. The map $\varphi:N(H)\to G(\tilde X)$ sending $[\gamma]$ to the
corresponding deck transformation is a surjective homomorphism (composition of
loops lifts to composition of the corresponding path then its deck-translate) with
kernel exactly the classes lifting to loops, i.e.\ $H$.
\end{proof}

\begin{definition}[Group action on a space; covering space action; free action]
\label{def:group-action}
\group{Covering Spaces}
An \textbf{action} of $G$ on $Y$ is a homomorphism $G \to \mathrm{Homeo}(Y)$. The
action is a \textbf{covering space action} if condition $(*)$ holds: each $y\in Y$ has a
neighborhood $U$ with $g(U)\cap U \ne \emptyset$ only for $g=$ identity (equivalently,
all $g(U)$ pairwise disjoint). The action is \textbf{free} if only the identity fixes any
point of $Y$; $(*)$ implies freeness but not conversely (e.g.\ irrational rotation of
$S^1$ is free but not a covering space action, since its orbits are dense, not
discrete).
\end{definition}

The deck transformation action of $G(\tilde X)$ on $\tilde X$ satisfies $(*)$: if
$g_1(\tilde U)\cap g_2(\tilde U)\ne\emptyset$ for $\tilde U$ mapped homeomorphically to some
$U\subset X$, then $g_1,g_2$ agree at a point of $\tilde U$ (since $\tilde U$ meets each
fiber at most once), forcing $g_1=g_2$.

\begin{definition}[Orbit space]
\label{def:orbit-space}
\group{Covering Spaces}
\uses{def:group-action}
For an action of $G$ on $Y$, the \textbf{orbit space} $Y/G$ has points the orbits
$Gy = \{g(y):g\in G\}$. For a normal covering space $\tilde X\to X$, $\tilde X/G(\tilde X) = X$.
\end{definition}

\begin{proposition}[Covering space actions give normal covering spaces]
\label{prop:covering-space-action-quotient}
\dcref{ch1:1.40}
\group{Covering Spaces}
\uses{def:group-action,def:orbit-space,prop:normal-covering-deck-group,prop:covering-unique-lifting}
If a $G$-action on $Y$ satisfies $(*)$, then: (a) $p:Y\to Y/G$ is a normal covering
space; (b) $G(\tilde X) \cong G$ if $Y$ is path-connected (identifying $\tilde X = Y$); (c) if
also $Y$ is locally path-connected, $G \cong \pi_1(Y/G)/p_*(\pi_1(Y))$.
\end{proposition}

\begin{proof}
(a) The sets $g(U)$ for $U$ as in $(*)$ are disjoint and each maps homeomorphically
to the open set $p(U)$ in the quotient topology, giving a covering space, normal since
$g_1g_2^{-1}$ carries $g_2(U)$ to $g_1(U)$. (b) $G$ injects into $G(\tilde X)$ and surjects
onto it since a deck transformation $f$ agrees with some $g\in G$ at any chosen
point (as $y,f(y)$ lie in the same orbit), hence $f=g$ by uniqueness of lifts. (c) is
\cref{prop:normal-covering-deck-group}(b) applied with $H=p_*\pi_1(Y)$, trivial in $\pi_1(Y)$
since $Y$ need not be simply-connected here -- more precisely this is the special
case $G(\tilde X)\cong\pi_1(X)/H$ with $\tilde X = Y$, $X = Y/G$.
\end{proof}

In particular, for a covering space action of $G$ on simply-connected, locally
path-connected $Y$, $\pi_1(Y/G) \cong G$ -- generalizing the computation of
$\pi_1(S^1)$ via $\Z$ acting on $\R$ by translation.

\begin{example}[The 11-hole torus covers a genus-3 surface]
\label{ex:genus-11-torus-covering}
\dcref{ch1:1.41}
\group{Covering Spaces}
\uses{prop:covering-space-action-quotient}
The closed orientable surface $M_{11}$ has a 5-fold rotational symmetry generating
a covering space action of $\Z_5$, with quotient $M_{11}/\Z_5 \cong M_3$: so $\pi_1(M_3)$
contains $\pi_1(M_{11})$ as a normal subgroup of index 5, quotient $\Z_5$. Replacing 2
holes per arm by $m$ and the 5-fold symmetry by $n$-fold gives
$M_{mn+1} \to M_{m+1}$; an exercise in Section 2.2 shows via Euler
characteristic that any covering $M_g\to M_h$ has $g=mn+1$, $h=m+1$ for some
$m,n$.


\end{example}

\begin{example}[Torus and Klein bottle as symmetry-group quotients of the plane]
\label{ex:grid-symmetries-torus-klein-bottle}
\dcref{ch1:1.42}
\group{Covering Spaces}
\uses{prop:covering-space-action-quotient}
Decorate the integer grid in $\R^2$ with arrows two ways: (i) all horizontal arrows
rightward, symmetry group $\Z\times\Z$ (all integer translations); (ii) horizontal
arrows alternating direction by row, symmetry group generated by translations
$(x,y)\mapsto(x+m,y+2n)$ together with glide reflections. Both actions satisfy
$(*)$ (minimum displacement 1), so $\pi_1(\text{torus})$ and $\pi_1(\text{Klein bottle})$
are these two symmetry groups, and the index-2 translation subgroup in case (ii)
realizes the torus as a double cover of the Klein bottle.
\end{example}

\begin{example}[$\pi_1(\RP^n) \cong \Z_2$ for $n\ge 2$]
\label{ex:real-projective-space-covering}
\dcref{ch1:1.43}
\group{Covering Spaces}
\uses{prop:covering-space-action-quotient,prop:pi1-sphere-trivial-n-geq-2,ex:real-projective-space-cw}
The antipodal map generates a $\Z_2$-action on $S^n$ with orbit space $\RP^n$
(\cref{ex:real-projective-space-cw}), a covering space action since antipodal open
hemispheres are disjoint. Since $S^n$ is simply-connected for $n\ge 2$
(\cref{prop:pi1-sphere-trivial-n-geq-2}), $\pi_1(\RP^n)\cong\Z_2$. A generator is the
projection of a path between antipodal points, and has order exactly 2 since its
square lifts to a loop in $S^n$, nullhomotopic there.
\end{example}

One may ask which other finite groups act freely on $S^n$: for even $n$ only $\Z_2$
(\cref{prop:z2-acts-freely-even-spheres}); for odd $n$ the cyclic groups $\Z_m$ always
act freely (via $v\mapsto e^{2\pi i\ell/m}v$ on $S^{2k-1}\subset\C^k$, orbit spaces the
\emph{lens spaces} of \cref{ex:lens-spaces}). There are also noncyclic
finite groups acting freely as rotations of $S^n$ for odd $n>1$, classified explicitly
in \cite{Wolf1984}; among noncyclic
groups, subgroups of the unit quaternions $S^3\subset\mathbb H$ act freely by
left-multiplication -- the generalized quaternion (binary dihedral) groups
$Q_{4m}=D^*_{4m}$, and the binary tetrahedral/octahedral/icosahedral groups
$T^*_{24},O^*_{48},I^*_{120}$ (double covers of the rotation groups of the regular
tetrahedron, octahedron/cube, icosahedron/dodecahedron via $S^3\to SO(3)$,
$u\mapsto(v\mapsto u^{-1}vu)$). A finite group acting freely on some $S^n$ must have
all abelian subgroups cyclic (no $\Z_p\times\Z_p$, sketched in an exercise for §4.2, not
yet transcribed) and at most one element of order 2 (\cite{Milnor1957}); groups with
cyclic abelian subgroups are completely classified (\cite{Brown1982}, section VI.9) --
the dihedral group $D_{2m}$ for odd $m>1$ satisfies this condition but has more than
one element of order 2. Conversely, both conditions together are also sufficient for
acting freely on \emph{some} sphere (\cite{MadsenThomasWall1976}, \cite{DavisMilgram1985}).

\begin{example}[The action of $G_{m,n}$ on the tree $T_{m,n}$]
\label{ex:gmn-action-on-tree}
\dcref{ch1:1.44}
\group{Covering Spaces}
\uses{ex:xmn-universal-cover-construction,ex:torus-knot-complement-fundamental-group}
Continuing \cref{ex:xmn-universal-cover-construction}, $G_{m,n} = \pi_1(X_{m,n})$
(presentation $\langle a,b\mid a^mb^{-n}\rangle$ from \cref{ex:torus-knot-complement-fundamental-group})
acts on $\tilde X_{m,n} = T_{m,n}\times\R$ preserving the product structure, hence
acting diagonally via actions on each factor: on $\R$, $a$ and $b$ act by
translation by $1/m$ and $1/n$ (unit-height rectangles), generating a cyclic
translation group of step $1/\mathrm{lcm}(m,n)$; on $T_{m,n}$, the action has kernel
the center $\langle a^m\rangle = \langle b^n\rangle$ (\cref{ex:torus-knot-complement-fundamental-group}),
inducing a non-free action of $\Z_m*\Z_n$ (elements $a,b$ and conjugates fix
vertices). Restricting to the kernel $K$ of $G_{m,n}\to\Z$ (the action on the $\R$
factor) gives a free, hence covering-space, action of $K$ on $T_{m,n}$, so $K$ is a
free group, $\pi_1(T_{m,n}/K)$.
\end{example}

\subsection*{Cayley Complexes}

\begin{definition}[Cayley graph and Cayley complex]
\label{def:cayley-graph-and-complex}
\group{Covering Spaces}
\uses{cor:every-group-realized-as-pi1}
For a presentation $G = \langle g_\alpha \mid r_\beta\rangle$, the \textbf{Cayley graph} has
vertex set $G$ and an edge from $g$ to $gg_\alpha$ for each generator $g_\alpha$; it is
connected since every element is a product of generators. Each relator $r_\beta$
determines a loop at every vertex $g$; attaching a 2-cell along each such loop gives
the \textbf{Cayley complex} $\tilde X_G$, on which $G$ acts by left multiplication (a
covering space action, matching the earlier $X_G$ of
\cref{cor:every-group-realized-as-pi1} as orbit space $\tilde X_G/G = X_G$).
\end{definition}

$\tilde X_G$ is in fact the universal cover of $X_G$: the homomorphism
$\varphi:\pi_1(X_G)\to G$ from \cref{prop:normal-covering-deck-group}'s proof sends the
edge class $[e_\alpha]$ to $g_\alpha$, hence is an isomorphism, so
$\ker\varphi = p_*(\pi_1(\tilde X_G)) = 0$, giving $\pi_1(\tilde X_G)=0$ since $p_*$ is
injective.

\begin{example}[Cayley complexes of small groups]
\label{ex:cayley-graph-free-group}
\dcref{ch1:1.45}
\group{Covering Spaces}
\uses{def:cayley-graph-and-complex}
For $G=\Z*\Z=\langle a,b\rangle$, $X_G = S^1\vee S^1$ and $\tilde X_G$ is the Cayley
graph of $\Z*\Z$ (the same universal cover pictured earlier), with $a,b$ acting as
shifts along the central horizontal/vertical axes.
\end{example}

\begin{example}[Cayley complex of $\Z\times\Z$]
\label{ex:cayley-complex-torus}
\dcref{ch1:1.46}
\group{Covering Spaces}
\uses{def:cayley-graph-and-complex}
For $G=\Z\times\Z=\langle x,y\mid xyx^{-1}y^{-1}\rangle$, $X_G$ is the torus and
$\tilde X_G = \R^2$ with vertex set $\Z^2$, $G$ acting by integer translation.
\end{example}

\begin{example}[Cayley complex of $\Z_n$]
\label{ex:cayley-complex-cyclic-group}
\dcref{ch1:1.47}
\group{Covering Spaces}
\uses{def:cayley-graph-and-complex,ex:real-projective-space-covering}
For $G=\Z_n=\langle x\mid x^n\rangle$, $X_G$ is $S^1$ with a disk attached via
$z\mapsto z^n$, and $\tilde X_G$ is $n$ disks with a common boundary circle, a
generator acting by cyclically permuting the disks via $2\pi/n$ rotation of the shared
boundary. ($n=2$ recovers $X_G=\RP^2$, $\tilde X_G=S^2$, cf.\
\cref{ex:real-projective-space-covering}.)
\end{example}

\begin{example}[Cayley complex of $\Z_2*\Z_2$]
\label{ex:cayley-complex-infinite-dihedral}
\dcref{ch1:1.48}
\group{Covering Spaces}
\uses{def:cayley-graph-and-complex,def:free-product-of-groups,ex:cayley-complex-cyclic-group}
For $G=\Z_2*\Z_2=\langle a,b\mid a^2,b^2\rangle$ (the infinite dihedral group, cf.\
\cref{def:free-product-of-groups}), the Cayley graph is an infinite chain of circles each
tangent to its neighbors; $\tilde X_G$ makes each circle the equator of a 2-sphere, an
infinite chain of tangent spheres, with the index-2 subgroup generated by $ab$
acting by even translations along the chain and the remaining elements acting as
antipodal maps composed with a flip of the whole chain. The orbit space is
$\RP^2\vee\RP^2$. The generalization to $\Z_m*\Z_n=\langle a,b\mid a^m,b^n\rangle$
arranges copies of the $\Z_m$ and $\Z_n$ Cayley complexes
(\cref{ex:cayley-complex-cyclic-group}) in a tree-like pattern.


\end{example}

\subsection*{Exercises}
% Exercises are transcribed for completeness but are not given blueprint labels or
% \uses edges.

\begin{enumerate}
\item For a covering space $p:\tilde X\to X$ and $A\subset X$, let $\tilde A=p^{-1}(A)$.
Show $p:\tilde A\to A$ is a covering space.

\item Show that if $p_1:\tilde X_1\to X_1$ and $p_2:\tilde X_2\to X_2$ are covering
spaces, so is $p_1\times p_2 : \tilde X_1\times\tilde X_2 \to X_1\times X_2$.

\item Let $p:\tilde X\to X$ have $p^{-1}(x)$ finite and nonempty for all $x$. Show
$\tilde X$ is compact Hausdorff iff $X$ is.

\item Construct a simply-connected covering space of $X\subset\R^3$, the union of
a sphere and a diameter; do the same for a sphere and a circle meeting it in two
points.

\item Let $X\subset\R^2$ be the boundary of the unit square together with the
vertical segments $x=1/2,1/3,\ldots$ inside it. Show every covering space of $X$
restricts to a homeomorphism over some neighborhood of the left edge, and deduce
$X$ has no simply-connected covering space.

\item For $X$ the shrinking wedge of circles (\cref{ex:shrinking-wedge-circles}) and
$\tilde X$ a certain covering space of it (a row of circles on a line), construct a
2-sheeted covering $Y\to\tilde X$ such that $Y\to\tilde X\to X$ is not a covering
space, though it still has unique path lifting.

\item Let $Y$ be a quasi-circle (part of the graph of $\sin(1/x)$, the segment
$[-1,1]$ on the $y$-axis, and a connecting arc). Collapsing the $y$-axis segment
gives $f:Y\to S^1$. Show $f$ does not lift to $\R\to S^1$ although $\pi_1(Y)=0$,
showing local path-connectedness is needed in the lifting criterion.

\item Let $\tilde X,\tilde Y$ be simply-connected covers of path-connected, locally
path-connected $X,Y$. Show $X\simeq Y \Rightarrow \tilde X\simeq\tilde Y$.

\item Show that if path-connected, locally path-connected $X$ has finite
$\pi_1(X)$, every map $X\to S^1$ is nullhomotopic.

\item Find all connected 2- and 3-sheeted covering spaces of $S^1\vee S^1$, up to
unbased isomorphism.

\item Construct finite graphs $X_1,X_2$ with a common finite-sheeted covering
space but no common covering target.

\item Draw the covering space of $S^1\vee S^1$ for the normal subgroup generated
by $a^2,b^2,(ab)^4$, and verify it.

\item Determine the 27-sheeted covering space of $S^1\vee S^1$ corresponding to
the subgroup generated by all cubes. [Related to Burnside's problem.]

\item Find all connected covering spaces of $\RP^2\vee\RP^2$.

\item For a simply-connected cover $p:\tilde X\to X$ and path-connected, locally
path-connected $A\subset X$ with $\tilde A$ a component of $p^{-1}(A)$, show
$p:\tilde A\to A$ corresponds to $\ker(\pi_1(A)\to\pi_1(X))$.

\item Given $X\to Y\to Z$ with $Y\to Z$ and $X\to Z$ covering spaces and $Z$
locally path-connected, show $X\to Y$ is a covering space, normal if $X\to Z$ is.

\item Given $G$ and normal $N\trianglelefteq G$, show there is a normal covering
space $\tilde X\to X$ with $\pi_1(X)\cong G$, $\pi_1(\tilde X)\cong N$,
$G(\tilde X)\cong G/N$.

\item Call a covering space \textbf{abelian} if normal with abelian deck group. Show
$X$ has a universal abelian covering space, unique up to isomorphism, covering
every other abelian covering space; describe it for $S^1\vee S^1$ and
$S^1\vee S^1\vee S^1$.

\item Show $M_g$ has a connected normal cover with deck group $\Z^n$ iff
$n\le 2g$; for $n=3$, $g\ge3$, realize it in $\R^3$ with translation deck
transformations iff $M_g$ embeds in the 3-torus $T^3$ with $\pi_1(M_g)\to\pi_1(T^3)$
surjective.

\item Construct nonnormal covers of the Klein bottle by a Klein bottle and by a
torus.

\item For $X$ the torus with a M\"obius band attached along $S^1\times\{x_0\}$,
compute $\pi_1(X)$ and describe the universal cover and the $\pi_1(X)$-action on
it; do the same for $Y$, $\RP^2$ with a M\"obius band attached along the 1-skeleton
circle.

\item Show the product action of $G_1\times G_2$ on $X_1\times X_2$ is a covering
space action if the factors are, with $(X_1\times X_2)/(G_1\times G_2) \cong
X_1/G_1\times X_2/G_2$.

\item Show a free, properly discontinuous action on a Hausdorff space is a
covering space action; in particular a free action of a finite group is.

\item For a covering space action of $G$ on path-connected, locally
path-connected $X$ and $H\subset G$: (a) every path-connected cover between $X$
and $X/G$ is $X/H$ for some $H$; (b) $X/H_1\cong X/H_2$ iff $H_1,H_2$ conjugate;
(c) $X/H\to X/G$ normal iff $H\trianglelefteq G$, with deck group $G/H$.

\item For $\varphi(x,y)=(2x,y/2)$ generating a $\Z$-action on
$X=\R^2-\{0\}$, show it is a covering space action, compute $\pi_1(X/\Z)$, and
show $X/\Z$ is non-Hausdorff, a union of four copies of $S^1\times\R$.

\item For $p:\tilde X\to X$ with $X$ connected, locally path-connected, semilocally
simply-connected: (a) components of $\tilde X$ correspond to orbits of
$\pi_1(X,x_0)$ on the fiber; (b) under the Galois correspondence, the subgroup for
the component containing $\tilde x_0$ is its stabilizer.

\item For the universal cover, show the fiber action via lifting loops and the
action via restricting deck transformations differ in general (e.g.\ for
$S^1\vee S^1$ and $S^1\times S^1$), and determine when they coincide.

\item Show $\pi_1(Y/G) \cong G$ for a covering space action on simply-connected
$Y$.

\item For $Y$ path-connected, locally path-connected, simply-connected and
$G_1,G_2\subset\mathrm{Homeo}(Y)$ defining covering space actions, show
$Y/G_1\cong Y/G_2$ iff $G_1,G_2$ are conjugate in $\mathrm{Homeo}(Y)$.

\item Draw the Cayley graph of $\Z\times\Z_2 = \langle a,b\mid b^2\rangle$.

\item Show the normal covers of $S^1\vee S^1$ are exactly the Cayley graphs of
2-generator groups (and similarly for $n$ circles and $n$-generator groups).

\item For covering spaces of CW complexes with cells projecting homeomorphically
(so restricting to a covering $\tilde X^1\to X^1$ of 1-skeleta): (a) two such covers
are isomorphic iff their restrictions to 1-skeleta are; (b) normality is detected on
1-skeleta; (c) the deck groups agree via restriction.

\item In \cref{ex:gmn-action-on-tree}, with $d=\gcd(m,n)$, $m'=m/d$, $n'=n/d$,
show $T_{m,n}/K$ has $m'$ vertices labeled $a$, $n'$ labeled $b$, and $d$ edges
between each $a$-$b$ pair, so $K$ is free of rank $dm'n'-m'-n'+1$.
\end{enumerate}

\section*{Additional Topics}
\subsection{1.A Graphs and Free Groups}
\label{sec:graphs-and-free-groups}

These Additional Topics (1.A, 1.B) are not part of the book's essential core, mainly
illustrating how covering space theory lets topology be used to study algebraic
properties of groups.

\begin{definition}[Graph, vertices, edges]
\label{def:graph}
\group{Fundamental Group}
\uses{def:cell-complex}
A \textbf{graph} is a 1-dimensional CW complex: a space $X$ built from a discrete
set $X^0$ (the \textbf{vertices}) by attaching 1-cells (the \textbf{edges}) $e_\alpha$, i.e.\ the quotient
of $X^0 \sqcup_\alpha I_\alpha$ identifying the endpoints of each $I_\alpha$ with points of $X^0$; an
edge does not include its endpoints, and $\bar e_\alpha \cong I$ or $S^1$. $X$ has the
\textbf{weak topology} with respect to the closures $\bar e_\alpha$, and is locally path-connected
(hence connected iff path-connected).
\end{definition}

A compact subset of a graph meets only finitely many vertices and edges (via a
discreteness argument on a transversal set $D$).

\begin{definition}[Subgraph and tree]
\label{def:subgraph-and-tree}
\group{Fundamental Group}
\uses{def:graph,def:contractible}
A \textbf{subgraph} of $X$ is a closed subspace that is a union of vertices and edges.
A \textbf{tree} is a contractible graph; a tree in $X$ is \textbf{maximal} if it contains every
vertex of $X$.
\end{definition}

\begin{proposition}[Maximal trees exist]
\label{prop:maximal-tree-exists}
\dcref{ch1:1A.1}
\group{Fundamental Group}
\uses{def:subgraph-and-tree}
Every connected graph contains a maximal tree, and every tree in the graph
extends to a maximal tree.
\end{proposition}

\begin{proof}
Given a subgraph $X_0$, build $X_0=X_0\subset X_1\subset X_2\subset\cdots$ by adjoining
at each stage the closures of all edges with an endpoint in the previous stage; the
union is all of $X$ (open and closed, $X$ connected). Build $Y_0=X_0\subset Y_1\subset\cdots$
by adjoining, at each stage, one edge from $Y_i$ to each new vertex of
$X_{i+1}-X_i$; then $Y=\bigcup Y_i$ deformation retracts onto $X_0$ (concatenating the
retractions $Y_{i+1}\to Y_i$ over shrinking time intervals) and contains every vertex
of $X$. Taking $X_0$ any subtree (e.g.\ a point) gives the result.
\end{proof}

Given a maximal tree $T\subset X$ and basepoint $x_0\in T$, each edge $e_\alpha$ of $X-T$
determines a loop $f_\alpha$ at $x_0$ (through $T$ to one end of $e_\alpha$, across $e_\alpha$, back
through $T$), well-defined up to homotopy since $T$ is simply-connected.

\begin{proposition}[$\pi_1$ of a connected graph is free]
\label{prop:fundamental-group-of-graph-is-free}
\dcref{ch1:1A.2}
\group{Fundamental Group}
\uses{prop:maximal-tree-exists,ex:wedge-sum-fundamental-group,prop:quotient-by-contractible-hep-homotopy-equiv}
For a connected graph $X$ with maximal tree $T$, $\pi_1(X)$ is free with basis the
classes $[f_\alpha]$ for edges $e_\alpha$ of $X-T$.
\end{proposition}

\begin{proof}
$X\to X/T$ is a homotopy equivalence by \cref{prop:quotient-by-contractible-hep-homotopy-equiv}
(Proposition 0.17), and $X/T$ is a one-vertex graph, i.e.\ a wedge of circles, with free
$\pi_1$ on the images of the $f_\alpha$'s by \cref{ex:wedge-sum-fundamental-group}.
\end{proof}

In particular a maximal tree is maximal in the naive sense too (adjoining any edge
creates nontrivial $\pi_1$), and a graph is a tree iff simply-connected.

\begin{lemma}[Covering spaces of graphs are graphs]
\label{lem:covering-space-of-graph-is-graph}
\dcref{ch1:1A.3}
\group{Covering Spaces}
\uses{def:graph,prop:covering-homotopy-property}
Every covering space of a graph is a graph, with vertices/edges the lifts of the base
graph's vertices/edges.
\end{lemma}

\begin{proof}
Take $\tilde X^0 = p^{-1}(X^0)$; path lifting applied to each $I_\alpha\to X$ gives a unique
lift through each point of $p^{-1}(x)$, $x\in e_\alpha$, defining the edges; the resulting
topology agrees with the covering space topology since $p$ is a local homeomorphism.
\end{proof}

\begin{theorem}[Every subgroup of a free group is free]
\label{thm:subgroup-of-free-group-is-free}
\dcref{ch1:1A.4}
\group{Covering Spaces}
\uses{prop:covering-space-existence-for-subgroup,prop:covering-map-pi1-injective,lem:covering-space-of-graph-is-graph,prop:fundamental-group-of-graph-is-free}
Every subgroup of a free group is free.
\end{theorem}

\begin{proof}
Realize $F$ as $\pi_1(X)$ for $X$ a wedge of circles; for $G\subset F$,
\cref{prop:covering-space-existence-for-subgroup} gives a covering $\tilde X\to X$ with
$p_*\pi_1(\tilde X)=G$, so $\pi_1(\tilde X)\cong G$ by injectivity
(\cref{prop:covering-map-pi1-injective}); $\tilde X$ is a graph by
\cref{lem:covering-space-of-graph-is-graph}, so $G$ is free by
\cref{prop:fundamental-group-of-graph-is-free}.
\end{proof}

\begin{definition}[Edgepath]
\label{def:edgepath}
\group{Fundamental Group}
\uses{def:graph,prop:maximal-tree-exists,lem:covering-space-of-graph-is-graph,def:universal-cover}
An \textbf{edgepath} is a composition of finitely many paths, each a single edge
traversed monotonically. In a tree, height functions (from the construction of
\cref{prop:maximal-tree-exists} starting at a single vertex) show there is a unique
\textbf{nonbacktracking} edgepath between any two vertices, and a graph with no
circle subgraph is a tree (else a maximal tree $T\ne X$ plus an edge of $X-T$ gives a
circle). More generally, for any connected graph $X$, each homotopy class of paths
between two vertices contains a unique nonbacktracking edgepath, obtained by
lifting to the (tree) universal cover.
\end{definition}

\subsection{1.B $K(G,1)$ Spaces and Graphs of Groups}
\label{sec:kg1-spaces-graphs-of-groups}

\begin{definition}[$K(G,1)$ space]
\label{def:k-g-1-space}
\group{Fundamental Group}
\uses{def:universal-cover,def:contractible}
A path-connected space with fundamental group $G$ and contractible universal
cover is a \textbf{$K(G,1)$ space} (an Eilenberg--Mac Lane space; more general
$K(G,n)$'s are studied in §4.2, not yet transcribed).
\end{definition}

\begin{example}[Graphs, surfaces, and lens spaces as $K(G,1)$'s]
\label{ex:graphs-are-kg1}
\dcref{ch1:1B.5}
\group{Fundamental Group}
\uses{def:k-g-1-space,prop:fundamental-group-of-graph-is-free,lem:covering-space-of-graph-is-graph}
A connected graph is a $K(F,1)$ for $F$ free, its universal cover being a tree
(\cref{prop:fundamental-group-of-graph-is-free}). $\RP^\infty$ is a $K(\Z_2,1)$: its universal
cover $S^\infty$ is contractible via an explicit two-stage straight-line homotopy (shift
coordinates, then contract to a point, each stage avoiding the origin so it descends
to $S^\infty$). More generally $S^\infty/\Z_m$ (scalar multiplication by $m$-th roots of unity
on $S^\infty\subset\C^\infty$) is a $K(\Z_m,1)$, an infinite-dimensional \textbf{lens space}. Products
of $K(G,1)\times K(H,1)$ are $K(G\times H,1)$ (universal cover = product of universal
covers), giving $K(G,1)$'s for finitely generated abelian $G$ from tori and lens
spaces; e.g.\ $T^n$ is a $K(\Z^n,1)$.
\end{example}

\begin{example}[Knot complements in $S^3$ as $K(G,1)$'s]
\label{ex:knot-complement-in-s3-kg1}
\dcref{ch1:1B.6}
\group{Fundamental Group}
\uses{def:k-g-1-space,ex:xmn-universal-cover-construction}
For a nonempty closed connected $K\subset S^3$, $S^3-K$ is a $K(G,1)$ (a theorem
of 3-manifold theory); for $K$ a torus knot this follows already from
\cref{ex:xmn-universal-cover-construction}, since $S^3-K$ deformation retracts onto
$X_{m,n}$, whose universal cover is contractible.
\end{example}

\begin{example}[Milnor's construction: $EG$ and $BG$]
\label{ex:milnor-construction-eg-bg}
\dcref{ch1:1B.7}
\group{Fundamental Group}
\uses{def:k-g-1-space,def:group-action,prop:covering-space-action-quotient}
For any group $G$, let $EG$ be the $\Delta$-complex with $n$-simplices the ordered
tuples $[g_0,\ldots,g_n]$ of $G$-elements, faces given by deleting entries. $EG$ is
contractible via a straight-line homotopy sliding each point toward the vertex
$[e]$ (well-defined on faces, though not literally a deformation retraction onto
$[e]$ since $[e]$ itself moves around the loop $[e,e]$). $G$ acts on $EG$ by left
multiplication, freely and permuting simplices (only $e$ fixes a simplex), hence a
covering space action (\cref{def:group-action}), so $BG := EG/G$ is a $K(G,1)$ with
universal cover $EG$. $BG$ inherits a one-vertex $\Delta$-complex structure, with
simplices $[g_1|\cdots|g_n]$ (`bar notation') and boundary faces
$[g_2|\cdots|g_n]$, $[g_1|\cdots|g_{n-1}]$, and $[g_1|\cdots|g_ig_{i+1}|\cdots|g_n]$. This
construction is functorial ($f:G\to H$ induces $Bf:BG\to BH$) but $BG$ is always
infinite-dimensional and, for finite $G$, has infinitely many cells in each positive
dimension -- far less efficient than special-case constructions like $S^1$ for $K(\Z,1)$.
A different, generally also inefficient, construction starts from any 2-dimensional
complex $X_G$ with $\pi_1(X_G)\cong G$ (\cref{cor:every-group-realized-as-pi1}) and
attaches cells of dimension $\ge 3$ to kill the higher homotopy of the universal cover
without disturbing $\pi_1$; it is a curious fact that if $G$ has an element of finite
order, every $K(G,1)$ CW complex is necessarily infinite-dimensional (shown in
\cref{prop:fundamental-group-kg1-torsionfree}) -- so, e.g., the infinite-dimensional lens spaces
$K(\Z_m,1)$ above cannot be replaced by finite-dimensional complexes.
\end{example}

In spite of this latitude, $K(G,1)$'s have a homotopical uniqueness property that
accounts for much of their interest:

\begin{theorem}[Homotopy type of a $K(G,1)$ is unique]
\label{thm:kg1-unique-homotopy-type}
\dcref{ch1:1B.8}
\group{Fundamental Group}
\uses{def:k-g-1-space,prop:homomorphism-to-kg1-realized-by-map}
The homotopy type of a CW complex $K(G,1)$ is uniquely determined by $G$. (Since
homology and cohomology depend only on homotopy type, this makes such
invariants of a $K(G,1)$ into invariants of the group $G$ itself.)
\end{theorem}

\begin{proof}
This follows from \cref{prop:homomorphism-to-kg1-realized-by-map}: given CW
$K(G,1)$'s $X,Y$ and an isomorphism $\pi_1(X,x_0)\cong\pi_1(Y,y_0)$, the proposition
gives maps $f:(X,x_0)\to(Y,y_0)$ and $g:(Y,y_0)\to(X,x_0)$ inducing inverse
isomorphisms on $\pi_1$; then $fg,gf$ induce the identity on $\pi_1$, so by the
uniqueness half of the proposition (applied to $\mathrm{id}_X,\mathrm{id}_Y$) they are
homotopic to the identity, exhibiting $f,g$ as homotopy inverses.
\end{proof}

\begin{proposition}[Maps into a $K(G,1)$ realize any homomorphism]
\label{prop:homomorphism-to-kg1-realized-by-map}
\dcref{ch1:1B.9}
\group{Fundamental Group}
\uses{def:k-g-1-space,def:cell-complex,prop:maximal-tree-exists}
Let $X$ be a connected CW complex and $Y$ a $K(G,1)$. Every homomorphism
$\varphi:\pi_1(X,x_0)\to\pi_1(Y,y_0)$ is induced by a map $(X,x_0)\to(Y,y_0)$, unique up
to homotopy fixing $x_0$.
\end{proposition}

\begin{proof}
Suppose first $X$ has a single $0$-cell $x_0$. Set $f(x_0)=y_0$ and, for each $1$-cell
$e^1_\alpha$ (giving $[e^1_\alpha]\in\pi_1(X,x_0)$), let $f$ on $\bar e^1_\alpha$ represent
$\varphi([e^1_\alpha])$; then $f_*=\varphi$ on $\pi_1(X^1,x_0)$, which is generated by the
$[e^1_\alpha]$'s. To extend $f$ over a $2$-cell $e^2_\beta$ with attaching map $\psi_\beta$, note
$i_*([\psi_\beta])=0$ in $\pi_1(X,x_0)$ (since $e^2_\beta$ already nullhomotopes $\psi_\beta$ in
$X$), so $f\psi_\beta$ represents $\varphi i_*([\psi_\beta])=0$ in $\pi_1(Y,y_0)$ and is
nullhomotopic, letting $f$ extend over $e^2_\beta$. For cells $e^n_\gamma$, $n>2$, the
attaching map $\psi_\gamma:S^{n-1}\to X^{n-1}$ has $f\psi_\gamma$ automatically
nullhomotopic: it lifts to the universal cover of $Y$ (contractible by
\cref{def:k-g-1-space}), and a nullhomotopy of the lift descends. Uniqueness up to
homotopy is proved one dimension up: two maps $f_0,f_1$ inducing the same $\varphi$
agree up to homotopy on $X^1$ (same generator argument), and this homotopy
extends over each cell $e^n\times(0,1)$ of $X\times I$ exactly as above, since these
cells have dimension $n+1>2$.

If $X$ has more than one $0$-cell, fix a maximal tree $T\subset X$
(\cref{prop:maximal-tree-exists}) and set $f(T)=y_0$; each edge $e^1_\lambda$ of $X-T$
still gives $[e^1_\lambda]\in\pi_1(X,x_0)$, and the construction and extension proceed
exactly as before. For the homotopy between two realizations $f_0,f_1$ of the same
$\varphi$, first homotope $f_0|X^1,f_1|X^1$ (via a deformation retraction of $T$ onto
$x_0$, extended to a homotopy of $X^1$) to agree on $T$, then apply the
single-vertex argument.
\end{proof}

The first half of this proof applies equally to the 2-dimensional complexes $X_G$ of
\cref{cor:every-group-realized-as-pi1}: every homomorphism $G\to H$ is induced by
some map $X_G\to X_H$. There is no uniqueness statement for such a map, though,
and different presentations of $G$ can give $X_G$'s that are not homotopy equivalent.

\begin{definition}[Graph of groups; graph product]
\label{def:graph-of-groups}
\group{Fundamental Group}
\uses{ex:milnor-construction-eg-bg,def:mapping-cylinder}
A \textbf{graph of groups} assigns a group $G_v$ to each vertex $v$ of a connected
oriented graph $\Gamma$ and a homomorphism $\varphi_e:G_{\mathrm{tail}(e)}\to G_{\mathrm{head}(e)}$ to
each edge $e$. Building $K\Gamma$ by placing $K(G_v,1)$ at each vertex and gluing in
the mapping cylinder of a map realizing $\varphi_e$ along each edge $e$ (identifying its
ends with the two vertex spaces), $\pi_1(K\Gamma)$ is called the \textbf{graph product} of the
$G_v$'s (\emph{`the fundamental group of the graph of groups'} in the literature).
\end{definition}

If $\Gamma$ is a central vertex with trivial group radiating edges to outer vertices, van
Kampen's theorem gives $\pi_1(K\Gamma)$ the free product of the outer vertex groups --
motivating the name.

\begin{theorem}[Graphs of groups with injective edge maps are $K(G,1)$'s]
\label{thm:graph-of-groups-is-kg1}
\dcref{ch1:1B.11}
\group{Fundamental Group}
\uses{def:graph-of-groups,def:k-g-1-space,prop:covering-space-action-quotient,cor:hep-inclusion-homotopy-equiv-deformation-retract}
If all edge homomorphisms $\varphi_e$ are injective, then $K\Gamma$ is a $K(G,1)$ and each
inclusion $K(G_v,1)\hookrightarrow K\Gamma$ is injective on $\pi_1$.
\end{theorem}

\begin{proof}
For $p:\tilde X\to X$ the universal cover and $A\subset X$ connected, locally
path-connected with $\pi_1(A)\to\pi_1(X)$ injective, each component of $p^{-1}(A)$ is
simply-connected (its $\pi_1$ injects into $\pi_1(\tilde X)=0$), hence is a universal cover
of $A$. Applying this with $X=M_f$ a mapping cylinder of $f:A\to B$: the lift
$\tilde M_f\to M_f$ is itself a mapping cylinder of a map $p^{-1}(A)\to p^{-1}(B)$, so
$\tilde M_f$ deformation retracts onto $p^{-1}(B)$, making $p^{-1}(B)$ simply-connected
too; if further $A,B$ are $K(G,1)$'s and $f_*$ injective, each component $\tilde A$ of
$p^{-1}(A)$ is a universal cover of $A$, hence contractible, and $\tilde M_f$ deformation
retracts onto it (both contractible, inclusion a homotopy equivalence, and
$(\tilde M_f,\tilde A)$ has the HEP by \cref{cor:hep-inclusion-homotopy-equiv-deformation-retract}
and the mapping-cylinder-neighborhood argument).

Build the covering space $\tilde K\to K\Gamma$ as an increasing union
$\tilde K_1\subset\tilde K_2\subset\cdots$: $\tilde K_1$ is the universal cover of one mapping
cylinder $M_f$; $\tilde K_{n+1}$ is obtained from $\tilde K_n$ by attaching, to each copy of a
universal cover of a $K(G_v,1)$ just created, a universal cover of each mapping
cylinder of $K\Gamma$ meeting it. Each stage deformation retracts onto the previous
(by the preceding paragraph), so concatenating these retractions over shrinking
time intervals (finishing with a contraction of $\tilde K_1$) shows $\tilde K=\bigcup_n\tilde K_n$ is
contractible, hence the universal cover of $K\Gamma$, proving $K\Gamma$ is a $K(G,1)$.

For injectivity of $K(G_v,1)\to K\Gamma$ on $\pi_1$: if $\gamma$ in $K(G_v,1)$ is nullhomotopic
in $K\Gamma$, the lifting criterion gives a lift $\tilde\gamma$ into $\tilde K$, landing inside one copy
of the universal cover of $K(G_v,1)$, so $\tilde\gamma$ (hence $\gamma$) is already nullhomotopic
there.
\end{proof}

The mapping cylinders making up $\tilde K$ are arranged along a tree $T_\Gamma$ (one vertex
per copy of a universal cover of some $K(G_v,1)$, edges from lifted mapping-cylinder
segments), built up inductively in step with the $\tilde K_n$; $\pi_1(K\Gamma)$ acts on $\tilde K$ by
deck transformations, inducing an action on $T_\Gamma$ with orbit space $\Gamma$, generally
not free (elements of $G_v$ fix the vertex for that copy of $K(G_v,1)$'s universal cover).
This exact correspondence between graphs of groups and group actions on trees is
developed further in \cite{ScottWall1979}.

\begin{example}[Free products with amalgamation]
\label{ex:free-product-with-amalgamation}
\dcref{ch1:1B.12}
\group{Fundamental Group}
\uses{thm:graph-of-groups-is-kg1,ex:xmn-universal-cover-construction}
For $\Gamma = A \leftarrow C \to B$ with both maps monomorphisms (embedding $C$ in
$A$ and $B$), van Kampen's theorem applied to the two mapping cylinders of $K\Gamma$
gives $\pi_1(K\Gamma) = A *_C B$, the free product of $A,B$ \textbf{amalgamated} along $C$,
containing both $A$ and $B$ as subgroups by \cref{thm:graph-of-groups-is-kg1}. Taking
$A=B=\Z$, $C=\Z$ via degree-$m$ and degree-$n$ covering maps $S^1\to S^1$
recovers $\Z *_\Z \Z$, realized (for coprime $m,n$) by $X_{m,n}$
(\cref{ex:xmn-universal-cover-construction}) as a deformation retract of a torus knot
complement -- and it is a theorem of 3-manifold theory that every smooth knot
complement in $S^3$ arises from iterated such constructions with injective edge maps
starting from free groups, hence is a $K(G,1)$ with universal cover $\R^3$.
\end{example}

\begin{example}[HNN extensions]
\label{ex:hnn-extension}
\dcref{ch1:1B.13}
\group{Fundamental Group}
\uses{thm:graph-of-groups-is-kg1}
For a graph of groups with a single vertex $A$ and a single edge with two
monomorphisms $\varphi,\psi:C\to A$ (a loop, i.e.\ $\Gamma$ has one vertex and one edge
mapping to itself), $\pi_1(K\Gamma)$, denoted $A*_C$, is a quotient of $A*\Z$ by relations
$t\varphi(c)t^{-1}=\psi(c)$ for $c\in C$ ($t$ a generator of the extra $\Z$); viewing $K\Gamma$ as
$K(A,1)\vee S^1$ with cells of dimension $\ge 2$ attached realizes this. Taking
$\varphi=\psi=\mathrm{id}$ gives $A*_A = A\times\Z$; taking $\varphi=\mathrm{id}$ and $\psi$ an arbitrary
automorphism realizes any semidirect product $A\rtimes\Z$ (e.g.\ the Klein bottle,
$\psi$ a reflection). A further example: deleting a disk from a torus and gluing the
resulting boundary circle to a longitude gives a space $X$ (homeomorphic to the
standard Klein-bottle-in-$\R^3$ picture) with
$\pi_1(X) = \langle a,b,t \mid tbt^{-1}aba^{-1}b^{-1}\rangle \cong (\Z*\Z)*_\Z\Z$; this abelianizes to
$\Z\times\Z$ but is \emph{not} free: setting $b=1$ gives a surjection
$\pi_1(X)\to\Z*\Z$ with nontrivial kernel (since $b\ne 1$ in $\pi_1(X)$ by
\cref{thm:graph-of-groups-is-kg1}), and if $\pi_1(X)\cong\Z*\Z$ this would give a
non-injective surjection $\Z*\Z\to\Z*\Z$, contradicting the fact (from group theory)
that free groups are \textbf{hopfian}.
\end{example}

\begin{example}[Closed surfaces are $K(G,1)$'s]
\label{ex:closed-surfaces-are-kg1}
\dcref{ch1:1B.14}
\group{Fundamental Group}
\uses{thm:graph-of-groups-is-kg1,ex:graphs-are-kg1}
A closed orientable surface $M$ of genus $\ge 2$ cuts along a circle into two
compact pieces $M_1,M_2$, each the mapping cylinder of the (non-nullhomotopic,
hence $\pi_1$-injective, since free groups are torsion-free) attaching map of the
2-cell in the standard CW structure applied to a punctured lower-genus closed
surface. This realizes $M = K\Gamma$ for $\Gamma$ of the form $F_1 \leftarrow \Z \to F_2$ with
$F_1,F_2$ free and both maps injective, so \cref{thm:graph-of-groups-is-kg1} shows $M$
is a $K(G,1)$ -- and, combined with the classification of simply-connected surfaces
($S^2$ or $\R^2$), $M$'s universal cover must be $\R^2$ since it is noncompact. A similar
argument (via M\"obius bands as mapping cylinders of the double cover $S^1\to S^1$)
handles closed nonorientable surfaces other than $\RP^2$. Non-closed surfaces
deformation retract onto graphs, hence are $K(F,1)$'s for $F$ free
(\cref{ex:graphs-are-kg1}).
\end{example}

\subsection*{Exercises}
% Exercises are transcribed for completeness but are not given blueprint labels or
% \uses edges.

\begin{enumerate}
\item Let $X$ be a graph in which each vertex has only finitely many incident
edges. Show the weak topology on $X$ is a metric topology.

\item Show a connected graph retracts onto any connected subgraph.

\item For a finite graph $X$, let $\chi(X) = |\text{vertices}|-|\text{edges}|$. Show
$\chi(X)=1$ if $X$ is a tree, and the rank of $\pi_1(X)$ is $1-\chi(X)$ if $X$ is connected.

\item If $X$ is a finite graph and $Y\subset X$ a subgraph homeomorphic to $S^1$
containing the basepoint, show $\pi_1(X,x_0)$ has a basis with one element
represented by $Y$.

\item Construct a connected graph $X$ and maps $f,g:X\to X$ with $fg=\one$ but
neither inducing an isomorphism on $\pi_1$.

\item For $F$ free on two generators with commutator subgroup $F'$, find free
generators for $F'$ via the covering space of $S^1\vee S^1$ corresponding to $F'$.

\item If $F$ is finitely generated free and $N\trianglelefteq F$ is nontrivial of infinite
index, show $N$ is not finitely generated.

\item Show a finitely generated group has finitely many subgroups of each given
finite index. [Do free groups first, via covering spaces of graphs.]

\item Show an index-$n$ subgroup $H\subset G$ has at most $n$ conjugates, and
deduce there is a normal $K\subset H$ of finite index in $G$.

\item For $X$ a wedge of $n$ circles and $\tilde X\to X$ a covering with finite connected
subgraph $Y\subset\tilde X$, show there is a finite graph $Z\supset Y$ on the same vertices
with $Y\to X$ extending to a covering $Z\to X$.

\item Combine the two preceding exercises to show finitely generated free groups
are \textbf{residually finite}: for $x\ne e$ in $F$ there is a finite-index normal
$H\not\ni x$.

\item For $F$ finitely generated free, $H\subset F$ finitely generated, $x\in F-H$, show
there is a finite-index $K\supset H$ with $x\notin K$.

\item For $x\ne e$ in a finitely generated free group $F$, show there is a finite-index
$H\subset F$ with $x$ part of a basis of $H$.

\item Show the existence of maximal trees is equivalent to the Axiom of Choice.

\item Suppose $G$ acts simplicially and freely on a $\Delta$-complex $X$. Show the
action is a covering space action.

\item For connected CW $X$ and $G$ with every homomorphism $\pi_1(X)\to G$
trivial, show every map $X\to K(G,1)$ is nullhomotopic.

\item Show every graph product of trivial groups is free.

\item Use van Kampen's theorem to compute $A*_C$ as a quotient of $A*\Z$, as
stated in the text.

\item For the graph of groups with one vertex $\Z$ and one edge realizing
multiplication by 2 (via the 2-sheeted cover $S^1\to S^1$), show $\pi_1(K\Gamma)$ has
presentation $\langle a,b\mid bab^{-1}a^{-2}\rangle$ (the first Baumslag--Solitar group) and
describe its universal cover as $T\times\R$.

\item Show that for a graph of groups with all edge maps injective $\Z\to\Z$, $K\Gamma$
can be given universal cover $T\times\R$; work out the infinite sequence
$\Z\xrightarrow{2}\Z\xrightarrow{3}\Z\xrightarrow{4}\Z\to\cdots$, showing $\pi_1(K\Gamma)\cong\Q$, and
describe a variant giving the dyadic rationals.

\item Show every graph product of groups is realized by a bipartite graph.

\item Show a finite graph product of finitely generated (resp.\ finitely presented)
groups is finitely generated (resp.\ finitely presented).

\item For a finite graph of finite groups with injective edge maps, show the graph
product has a free subgroup of finite index. [The converse also holds: see
\cite{ScottWall1979}, Theorem 7.3.]
\end{enumerate}
